%%%%%%%% mlsys 2025 EXAMPLE LATEX SUBMISSION FILE %%%%%%%%%%%%%%%%%

\makeatletter
\providecommand\Notice@String[1]{#1}
\makeatother

\documentclass{article}

% Recommended, but optional, packages for figures and better typesetting:
\usepackage{microtype}
\usepackage{graphicx}
% \usepackage{subfigure}
\usepackage{subcaption}
\usepackage{booktabs} % for professional tables

% hyperref makes hyperlinks in the resulting PDF.
% If your build breaks (sometimes temporarily if a hyperlink spans a page)
% please comment out the following usepackage line and replace
% \usepackage{mlsys2025} with \usepackage[nohyperref]{mlsys2025} above.
\usepackage{hyperref}

% Attempt to make hyperref and algorithmic work together better:
\newcommand{\theHalgorithm}{\arabic{algorithm}}

% Use the following line for the initial blind version submitted for review:
\usepackage[accepted]{mlsys2025}

\usepackage{enumitem}
\usepackage{booktabs}
\usepackage{inconsolata}
\usepackage{subcaption}
\usepackage{amsmath}
\usepackage{listings}

\definecolor{codebg}{RGB}{248,248,248}
\definecolor{coderule}{RGB}{215,215,215}
\lstdefinestyle{codebox}{
  basicstyle=\ttfamily\scriptsize,
  backgroundcolor=\color{codebg},
  frame=single,
  rulecolor=\color{coderule},
  framerule=0.4pt,
  columns=fullflexible,
  keepspaces=true,
  breaklines=true,
  breakatwhitespace=false,
  showstringspaces=false,
  tabsize=4,
  aboveskip=0.3em,
  belowskip=0.0em,
  xleftmargin=0.35em,
  xrightmargin=0.15em,
}

\usepackage{quoting}

\newcommand{\lily}[1]{\textcolor{blue}{Lily: #1}}
\newcommand{\jerry}[1]{\textcolor{cyan}{Jerry: #1}}
\newcommand{\alvin}[1]{\textcolor{red}{alvin: #1}}
\newcommand{\jong}[1]{\textcolor{purple}{jong: #1}}
\newcommand{\ion}[1]{\textcolor{blue}{ion: #1}}

\renewcommand{\subsectionautorefname}{Sec.}
\renewcommand{\sectionautorefname}{Sec.}
\renewcommand{\figureautorefname}{Fig.}
\renewcommand{\tableautorefname}{Tab.}

% If accepted, instead use the following line for the camera-ready submission:
% \usepackage[accepted]{mlsys2025}

% The \mlsystitle you define below is probably too long as a header.
% Therefore, a short form for the running title is supplied here:
% \mlsystitlerunning{Submission and Formatting Instructions for MLSys 2025}
\mlsystitlerunning{Speculative Decoding: Performance or Illusion?}

\begin{document}

\twocolumn[
\mlsystitle{Speculative Decoding: Performance or Illusion?}

% It is OKAY to include author information, even for blind
% submissions: the style file will automatically remove it for you
% unless you've provided the [accepted] option to the mlsys2025
% package.

% List of affiliations: The first argument should be a (short)
% identifier you will use later to specify author affiliations
% Academic affiliations should list Department, University, City, Region, Country
% Industry affiliations should list Company, City, Region, Country

% You can specify symbols, otherwise they are numbered in order.
% Ideally, you should not use this facility. Affiliations will be numbered
% in order of appearance and this is the preferred way.
\mlsyssetsymbol{equal}{*}

\begin{mlsysauthorlist}
\mlsysauthor{Xiaoxuan Liu}{equal,to}
\mlsysauthor{Jiaxiang Yu}{equal,to}
\mlsysauthor{Jongseok Park}{to}
\mlsysauthor{Ion Stoica}{to}
\mlsysauthor{Alvin Cheung}{to}
% \mlsysauthor{Tateu H.~Yasehe}{ed,to,goo}
% \mlsysauthor{Aaoeu Iasoh}{goo}
% \mlsysauthor{Buiui Eueu}{ed}
% \mlsysauthor{Aeuia Zzzz}{ed}
% \mlsysauthor{Bieea C.~Yyyy}{to,goo}
% \mlsysauthor{Teoau Xxxx}{ed}
% \mlsysauthor{Eee Pppp}{ed}
\end{mlsysauthorlist}

\mlsysaffiliation{to}{UC Berkeley}
% \mlsysaffiliation{goo}{Googol ShallowMind, New London, Michigan, USA}
% \mlsysaffiliation{ed}{School of Computation, University of Edenborrow, Edenborrow, United Kingdom}

\mlsyscorrespondingauthor{Xiaoxuan Liu}{xiaoxuan\_liu@berkeley.edu}
\mlsyscorrespondingauthor{Jiaxiang Yu}{jxyu@berkeley.edu}

% You may provide any keywords that you
% find helpful for describing your paper; these are used to populate
% the "keywords" metadata in the PDF but will not be shown in the document
\mlsyskeywords{Machine Learning, MLSys}

\vskip 0.3in

\begin{abstract}
Speculative decoding (SD) has become a popular technique to accelerate Large Language Model (LLM) inference, yet its real-world effectiveness remains unclear as prior evaluations rely on research prototypes and unrealistically small batch sizes. We present, to our knowledge, the first systematic study of SD on a production-grade and widely deployed inference engine (vLLM), covering multiple SD variants ($n$-gram, EAGLE/EAGLE-3, Draft-Model, Multi-Token Prediction) across diverse workloads, model scales, and batch sizes. We analyze key factors governing SD performance, and 
quantify a theoretical upper bound on SD speedup. Our results show that verification by the target model dominates the execution, while acceptance length varies markedly across output token positions, requests, and datasets. Comparing measured performance with theoretical bounds reveals substantial gaps between observed and theoretical upper bounds, and we leverage this observation to highlight new research opportunities that our study opens up in improving SD. Our code for profiling and simulator is available at \href{https://github.com/orgs/SpecDecode-Bench/repositories}{https://github.com/orgs/SpecDecode-Bench/repositories}.
\end{abstract}
]

% this must go after the closing bracket ] following \twocolumn[ ...

% This command actually creates the footnote in the first column
% listing the affiliations and the copyright notice.
% The command takes one argument, which is text to display at the start of the footnote.
% The \mlsysEqualContribution command is standard text for equal contribution.
% Remove it (just {}) if you do not need this facility.

%\printAffiliationsAndNotice{}  % leave blank if no need to mention equal contribution
\printAffiliationsAndNotice{\mlsysEqualContribution} % otherwise use the standard text.



%% >>> begin intro.tex
\section{Introduction}
Speculative decoding (SD) has emerged as a highly effective approach to speed up inference in large language models (LLMs).
It has been widely used in real workloads, ranging from chat completion~\cite{chatgpt}, question answering~\cite{cobbe2021gsm8k, rein2023gpqagraduatelevelgoogleproofqa} to coding tasks~\cite{jimenezswe}. 

Despite the substantial progress in SD, there has not been a systematic study on its effectiveness.
Prior studies only used prototype implementations rather than production-level systems, which greatly undermines the validity of their conclusions. 
Worse yet, these evaluations are often conducted with a batch size of one~\cite{xia-etal-2024-unlocking, cai2024medusa, leviathan2023fast, li2024eagle, li2024eagle2fasterinferencelanguage, li2025eagle3}—an unrealistic configuration that fails to reflect real-world deployment.
Second, since the introduction of SD, numerous variants have emerged, including draft-model-based~\cite{miao2023specinfer, liu2023online, zhou2023distillspec, cascade-inference, chen2024sequoia}, $n$-gram-based~\cite{saxena2023prompt, somasundaram2024pld+}, and tree-based~\cite{cai2024medusa, li2024eagle, lin2024bita, fu2024break, li2024eagle2fasterinferencelanguage, li2025eagle3} approaches. Yet, we are unaware of any systematic analysis comparing these variants and understanding their respective use cases. This gap poses practical challenges: when deploying SD in production, practitioners are often left without clear guidance on which variant to use for different model architectures or workload conditions.

We aim to study SD in actual deployment settings using real-world workloads. Towards that end, we systematically benchmark SD in vLLM~\cite{kwon2023efficient}, a production-ready inference system, and evaluate multiple SD variants across diverse datasets. 
% To capture different system loads, we vary the batch size to simulate different system loads. 
Although this setup may appear straightforward, several subtle yet important caveats arise in practice. %\alvin{caveats or surprising conclusion?} \jerry{I think "caveats" are better.}
First, while SD is theoretically guaranteed to preserve the same token distribution as standard decoding, we observe that the generated outputs are not always identical. This discrepancy stems from the inherent nondeterminism~\cite{he2025nondeterminism} in LLM inference—caused by factors such as kernel parallelism and floating-point variation. 
% To ensure a fair comparison, we filter out-of-distribution samples \alvin{what are ood samples} and constrain the generation lengths to be within a comparable range \alvin{comparable with what?}, thereby ensuring the reported speedup metrics are meaningful and consistent.\jerry{We are currently not using this length-based filter, remove this sentence?}
Furthermore, we study the SD performance on two increasingly prominent setups: reasoning workloads, which involve longer and more structured generations, and multi-token prediction (MTP)~\cite{deepseekai2025deepseekv3technicalreport, 5team2025glm45agenticreasoningcoding}, a recently emerging acceleration technique that complements SD.
Across all workloads, every SD variant outperforms the no-SD baseline. The speedup decreases as the batch size increases, consistent with the reduced opportunities for SD on computation-bound scenarios~\cite{liu2024dsd}. Notably, reasoning workloads that require long chains of thought—and therefore generate longer outputs—exhibit speedups comparable to those observed on standard language-modeling tasks.

Next, we conduct a detailed analysis to understand the performance of speculative decoding (SD) across different variants. The overall speedup of SD is governed by two key factors: the execution efficiency of its constituent stages and the token acceptance rate.

We first examine the execution efficiency of the different SD stages.
In general, SD consists of three stages---proposing tokens, verifying tokens, and system overhead such as scheduling and rejection sampling. 
By examining the time breakdown of the three stages, we find that the proposing stage accounts for only a small portion of the total execution time, and the execution of the large model during verification remains the dominant cost. Thus, when the acceptance rate is low, repeatedly running the large model on rejected tokens incurs substantial runtime cost. As shown in prior work~\cite{liu2024dsd}, this extra cost becomes prohibitive under high system load, and can even cause SD to be slower than standard decoding.

We next conduct a detailed analysis of acceptance behavior in SD, examining when proposed tokens are accepted or rejected.
We observe that acceptance patterns exhibit three levels of variability: within a request, across different requests, and across different datasets. 
Furthermore, different SD methods demonstrate entirely different behavior. For example, learned SD methods such as EAGLE maintain stable and consistently high acceptance across diverse workloads, whereas training-free method like $n$-gram produces more variable but occasionally much longer accepted spans in tasks like code-editing, where frequent local repetitions allow it to reuse previously seen patterns.

Based on the above observations, we ask a fundamental question: \emph{how far are current SD approaches from the optimal speedup, and what is the theoretical upper bound of speculative decoding?}

First, motivated by the observation that a large fraction of proposed tokens are ultimately rejected, we identify a promising direction for improving SD performance: \textbf{verifying only those tokens that are likely to be accepted}. In the ideal case, a proposal method that perfectly predicts the large model’s outputs would eliminate the need to invoke the large model altogether. To quantify the gap between observed and theoretical speedups, we develop a simulator grounded in real-world benchmarking data. This simulator evaluates performance under an idealized setting in which all proposed tokens are accepted, thereby minimizing verification cost and achieving the optimal speedup.

Second, we observe that different SD methods achieve higher acceptance at different token positions. This complementarity suggests that adaptively combining multiple SD methods can unlock additional speedup beyond what any single method can achieve in isolation. Our results show that such adaptive combinations can improve end-to-end speedup to 4.9$\times$ relative to standard (no-SD) decoding.

In summary, the paper makes the following contributions:
\begin{itemize}[leftmargin=*, itemsep=2pt]
    \item \textbf{Production-grade evaluation of speculative decoding.} 
    We present the first systematic study of SD within a widely adopted, highly optimized inference engine (vLLM), thereby bridging the gap between research prototypes and real-world deployments. We evaluate mainstream SD variants across various real-world workloads.

    \item \textbf{Decomposition of speedup and performance bottlenecks.} 
    We analyze key factors that affect SD performance and find that: (1) verification cost remains dominant, and (2) acceptance behavior exhibits variability across positions, requests, and datasets. To the best of our knowledge, this is the first work to thoroughly examine both the time/memory breakdown and acceptance dynamics of SD across diverse workloads.

    \item \textbf{Theoretical upper bound of SD speedup.} 
    Based on these observations, we analyze the gap between measured and maximum achievable speedup. By leveraging position-specific acceptance statistics without modifying the proposing method, we compute the theoretical upper bound of SD performance, revealing an orthogonal and promising direction for future optimization.
\end{itemize}















%% <<< end intro.tex


%% >>> begin background.tex
\section{Background and Motivation}
\textbf{Speculative Decoding.}
Large language models (LLMs) generate text auto-regressively: given a prompt $(x_1, \ldots, x_n)$, they produce an output sequence $(x_{n+1}, \ldots, x_{n+S})$, 
one token at a time. At each step, the model computes a probability distribution over the vocabulary and samples the next token. 
% This strictly sequential decoding limits the utilization of highly parallel hardware such as GPUs~\cite{pope2023efficiently} and results in high latency. 
However, while LLMs \emph{generate} tokens sequentially, they can also \emph{evaluate} the probabilities of a \emph{given} token sequence in parallel. That is, for candidate tokens $x_{n+1}, \ldots, x_{n+S}$, an LLM can compute $P(x_{n+1} \mid x_1, \ldots, x_n), \ldots, P(x_{n+S} \mid x_1, \ldots, x_{n+S-1})$ in a single forward pass.
Speculative decoding~\cite{leviathan2023fast, chen2023accelerating} utilizes this property of LLMs as an evaluator on a string of candidate tokens in parallel, enabling higher hardware utilization and potentially generating multiple tokens.
For example, a smaller model may generate $k$ candidate tokens, and our target LLM model then predicts $k+1$ probability distributions for each token position in parallel. We then accept the first $m$ tokens from the $k$ candidate tokens that satisfy the sampling method. 

\textbf{SD Variations.}
When SD was first introduced, the standard approach was to use a smaller draft model with the same vocabulary as the target model, since rejection sampling must operate on a probability distribution in the same token embedding space.
Subsequent work has proposed various ways to construct the draft model, such as quantization~\cite{lin2023awq, frantar2023gptq, kim2024squeezellm} or distillation~\cite{zhou2023distillspec} of the target LLM.
However, in practice, an off-the-shelf draft model is not always available, requiring quantizing or pruning the original model, or even pretraining a smaller model from scratch.

To overcome the limitations of requiring a separate draft model, \emph{draft-model-free} SD methods have been proposed.
One common approach is to add auxiliary prediction layers on top of the final transformer block ~\cite{li2024eagle, li2024eagle2fasterinferencelanguage, li2025eagle3}.
These additional ``head'' layers take the hidden states as input and directly propose tokens.
However, practitioners still need to fine-tune these drafting heads separately.
% Another line of work fully eliminates the reuqirement of training. \emph{Multi-Token Prediction (MTP)}~\cite{deepseekai2025deepseekv3technicalreport, 5team2025glm45agenticreasoningcoding} co-trains the auxiliary heads jointly with the main model, rather than fine-tuning them afterwards. As a result, auxiliary heads are provided out-of-the-box with the model, while also achieving a significantly higher token acceptance rate compared to the post-training, add-on approaches.
Recently, a new line of work has removed the need for post-hoc training. Multi-Token Prediction (MTP)~\cite{deepseekai2025deepseekv3technicalreport,5team2025glm45agenticreasoningcoding} co-trains auxiliary prediction heads jointly with the main model, rather than fine-tuning them afterward. As a result, these auxiliary heads are available out of the box and achieve substantially higher token-acceptance rates.
% highlighting the potential of training-time co-design \alvin{what's training time co design?} for speculative decoding.

Another line of draft-model-free SD methods utilize a non-LLM proposer.
A representative example is prompt lookup decoding, or \emph{$n$-gram}~\cite{saxena2023prompt}.
The intuition is that LLM outputs often contain recurring phrases; therefore, we can extract $n$-gram snippets from previously generated text and reuse them as token proposals.
This approach is particularly effective in workloads such as code or passage editing, where the input and output share substantial overlap. 
Another example is using retrieval-augmented generation (RAG) to retrieve candidate proposals from an external corpus~\cite{he2024restretrievalbasedspeculativedecoding}.
These non-LLM approaches demonstrate the potential for lightweight or task-specific drafting mechanisms.
% However, it remains an open problem to determine how to construct the most effective non-LLM proposers, and how to combine them optimally with standard SD pipelines.

\textbf{Problems of Existing Benchmarks.}
Most prior works evaluate SD by building proof-of-concept research prototypes, as integration into a production-grade inference system is nontrivial and requires modifications across multiple components.
However, this makes a fair comparison between SD methods difficult and the benefits in real-world deployment impossible to predict, due to the following limitations.
First, many prototype implementations lack key optimizations in production systems, such as the use of CUDA graphs~\cite{cudagraph} to reduce system overhead, a common feature in LLM serving systems~\cite{kwon2023efficient, zheng2024sglang}.
Second, the implementations are typically evaluated with a batch size of one, which is unrepresentative of large-scale batched deployments of the real world.
Finally, prior work lacks an in-depth analysis of SD speedup, as they primarily focus on high-level metrics such as average latency or dataset-level acceptance rate. \cite{li2024eagle,li2025eagle3,xia-etal-2024-unlocking}
We believe that a deeper understanding, such as execution time breakdown and position-level acceptance behavior, is necessary to fully explain and optimize the end-to-end performance of SD.
%% <<< end background.tex


%% >>> begin performance.tex
\section{End-to-End Performance of Speculative Decoding}
\label{sect:e2e-perf}
In this section, we evaluate the performance of SD variants across various model sizes and datasets.
% The objective is to obtain a comprehensive view of its effectiveness in a real serving system.
The complete experimental setup is summarized in  \autoref{tab:exp-config} of the Appendix.
% \jerry{@Lily, I shifted it to Appendix.}

\subsection{Experiment Setup}
\textbf{Inference Engine and Hardware.}
All experiments are conducted using vLLM v0.10.1.1 unless otherwise stated.
To closely approximate real-world serving scenarios, we enable all default vLLM optimizations, 
including KV cache management, continuous batching, chunked prefill, and CUDA Graphs.
We run all experiments on NVIDIA H100~\cite{h100} with 80 GB of memory.
For 8B models, we use a single GPU, while for the 70B/106B models, we use four GPUs with a tensor parallel size of four. 
% \jerry{Citation}

\textbf{Models and SD variants}
We evaluate two models from the Llama-3 family \cite{dubey2024llama}:
Llama3.1-8B-Instruct and Llama3-70B-Instruct.
We also include Qwen3-8B\footnote{Throughout the paper, Qwen3-8B-Thinking denotes the same model with thinking mode enabled, whereas Qwen3-8B refers to its non-thinking mode.}~\cite{yang2025qwen3technicalreport} and GLM-4.5-Air-106B \cite{5team2025glm45agenticreasoningcoding} for reasoning workloads. We benchmark the following SD variants:
\begin{itemize} [leftmargin=1.5em, itemsep=0pt]
    \item \textbf{Draft-model-based}~\cite{leviathan2023fast, chen2023accelerating}: Uses a smaller draft model to propose 3 draft tokens per step, which are later verified by the target model. Details of model pairing are in Appendix~\autoref{appendix:model-config}.
  \item \textbf{EAGLE}~\cite{li2024eagle} and \textbf{EAGLE-3}~\cite{li2025eagle3}: Draft-model-free methods with fine-tuned auxiliary prediction heads, predicting 3 draft tokens per step for chain-based settings\footnote{Official EAGLE-3 weights are unavailable for Llama-3-70B, EAGLE weights are likewise unavailable for Qwen3-8B and EAGLE-3/EAGLE weights are also unavailable for GLM-4.5-Air~\cite{li2024eagle}.}.
  \item \textbf{Multi-Token Prediction} (MTP)~\cite{deepseekai2025deepseekv3technicalreport}: Co-trained draft-model-free approach where auxiliary heads are jointly trained with the target model; we use 3 draft tokens per step.
  \item \textbf{$n$-gram}~\cite{saxena2023prompt, somasundaram2024pld+}: Training-free method reusing recurring $n$-gram phrases from the prompt; we use 3 draft tokens per step with {\tt prompt\_lookup\_max=7} and {\tt prompt\_lookup\_min=3}. For \texttt{InstructCoder} with Llama3.1-8B, Llama3-70B and Qwen3-8B without reasoning, we additionally evaluate $n$-gram with 5 draft tokens per step to study the effect of a larger proposal length.
  
\end{itemize}
% \begin{itemize} [leftmargin=1.5em, itemsep=0pt]
%   \item \textbf{$n$-gram}~\cite{saxena2023prompt, somasundaram2024pld+}: 
%    5 draft tokens per step, with {\tt prompt\_lookup\_max=7} and {\tt prompt\_lookup\_min=3};
%   \item \textbf{EAGLE}~\cite{li2024eagle} and \textbf{EAGLE-3}~\cite{li2025eagle3} 3 draft tokens per step and the official draft models from the EAGLE authors.
%   \item \textbf{Draft-Model Based}~\cite{leviathan2023fast, chen2023accelerating}: 3 draft tokens per step.
%   \item \textbf{Multi-Token Prediction} (MTP)~\cite{deepseekai2025deepseekv3technicalreport}: 3 draft tokens per step. \jerry{Add specific eagle model used, smaller draft model used in appendix}
% \end{itemize}


Unless otherwise specified, the maximum generation length is capped at 8192 tokens.
For reasoning workloads AIME22-24 and GPQA-Main, we extend this limit to 32,768 tokens to prevent early truncation.
Decoding is performed with temperature set to 0 and no sampling truncation (\texttt{top\_p=1}, \texttt{top\_k=-1}).
% to minimize differences in output lengths across different methods.

\textbf{Workloads.}
We evaluate six datasets representing diverse real-world applications:
\textit{CNN/DailyMail}~\cite{hermann2015teaching, see-etal-2017-get} (summarization),
\textit{ShareGPT}~\cite{sharegpt} (multi-turn chat),
\textit{InstructCoder}~\cite{li2024instructcoder} (code editting),
\textit{GSM8K}~\cite{cobbe2021gsm8k} (grade-school math),
and two complex reasoning datasets: \textit{AIME22-24}~\cite{AIME-AoPS, AIME-HF} and \textit{GPQA-Main}~\cite{rein2023gpqagraduatelevelgoogleproofqa}. We detailed the dataset information in Appendix \autoref{appendix:dataset}. 
% \jerry{Add more description;}


% For each dataset, we randomly sample up to \textbf{500} requests (or the dataset maximum if smaller). To emulate different system loads, we vary the batch size from 1 to 128 requests during inference. Reasoning workloads use smaller batch sizes from 1 to 4 due to their longer generation lengths. Figure~\ref{fig:dataset} shows the input and output length distributions for all datasets using Qwen3-8B after filtering. \jerry{TODO: readers wont know what filtering means here}.

% \begin{figure}[t]
%     \centering
%     \begin{subfigure}[b]{\columnwidth}
%         \centering
%         \includegraphics[width=\linewidth]{figures/workload_length_distribution/after_filtering/input_length_distributions_Qwen3-8B.pdf}
%         \caption{Input length distributions.}
%         \label{fig:data-input}
%     \end{subfigure}
%     \begin{subfigure}[b]{\columnwidth}
%         \centering
%         \includegraphics[width=\linewidth]{figures/workload_length_distribution/after_filtering/output_length_distributions_Qwen3-8B.pdf}
%         \caption{Output length distributions.}
%         \label{fig:data-output}
%     \end{subfigure}

%     \caption{Dataset input and output length distributions. \lily{change tick to 5K, 10K...}}
%     \label{fig:dataset}
% \end{figure}
% \jerry{TODO: check with Lily on which figure to use for input/output length figure, add citations for datasets}


\textbf{Generation Length and Evaluation Metrics.}
Although in theory SD should preserve the same token distribution as standard decoding, in practice we observe discrepancies in the generated outputs—-{\em even under greedy decoding.} \autoref{tab:gen-len} illustrates this by showing variations in generation length when using the same set of prompts while changing the batch size and SD variant.
For each configuration, we replicate each request multiple times to match the target batch size.
We repeat the process for 500 requests and record the generation length for each request.
\begin{table}[h]
\centering
\caption{Average generation length (mean $\pm$ standard deviation) of the baseline (w/o SD) and different SD variants on the ShareGPT dataset using Llama3.1-8B across varying batch sizes.}
\label{tab:gen-len}
\resizebox{0.9\linewidth}{!}{
\begin{tabular}{ccccc}
\hline
\textbf{Batch Size} & \textbf{w/o SD} & \textbf{ngram} & \textbf{Eagle} & \textbf{Eagle3} \\
\hline
1   & 737$\pm$1048 & 750$\pm$1100 & 731$\pm$1010 & 736$\pm$1046 \\
8   & 722$\pm$994  & 676$\pm$810  & 709$\pm$928  & 704$\pm$916  \\
16  & 708$\pm$932  & 675$\pm$818  & 686$\pm$853  & 686$\pm$853  \\
32  & 674$\pm$784  & 697$\pm$887  & 693$\pm$857  & 706$\pm$917  \\
64  & 687$\pm$852  & 696$\pm$911  & 709$\pm$924  & 709$\pm$924  \\
128 & 714$\pm$930  & 739$\pm$1050 & 709$\pm$924  & 709$\pm$924  \\
\hline
\end{tabular}
}
\end{table}

As shown in \autoref{tab:gen-len}, even without speculative decoding (the w/o SD column), generation length differs across different batch sizes, and similarly when SD is used. Importantly, these differences are not consistently longer or shorter than the baseline, suggesting that they stem from numerical nondeterminism rather than implementation bugs.
Recent work~\cite{he2025nondeterminism} attributes such variability to kernel-level nondeterminism and floating-point variation during GPU execution.

To account for fluctuations in output length, we use token \textbf{throughput}, i.e., the number of generated tokens per second, as the performance metric.
This measure removes the confounding effect of varying generation lengths. \textbf{Speedup} is then defined as the ratio between the throughput with speculative decoding (SD) and the throughput without it.

% Our primary evaluation metric is the \textbf{average end-to-end request latency}, measured from request arrival to completion.
% We compare this latency against standard (non-SD) decoding under identical configurations and report the resulting speedup ratio.
% Each configuration (model, dataset, SD variant, and batch size) is executed once after a short warmup phase.

\subsection{Results}
\begin{figure*}[h!]
\begin{center}
\captionsetup{justification=centering}
\includegraphics[width=0.6\linewidth]{figures/e2e/legend.pdf}
  
\begin{subfigure}[t]{0.22\textwidth}
\includegraphics[width=\linewidth]{figures/e2e/gen_speedup_llama3-8b_gsm8k.pdf}
\caption{Llama3.1-8B, GSM8K}
\end{subfigure}
\begin{subfigure}[t]{0.22\textwidth}
\includegraphics[width=\linewidth]{figures/e2e/gen_speedup_llama3-8b_cnndailymail.pdf}
\caption{Llama3.1-8B, CNN/Dailymail}
\end{subfigure}
\begin{subfigure}[t]{0.22\textwidth}
\includegraphics[width=\linewidth]{figures/e2e/gen_speedup_llama3-8b_sharegpt.pdf}
\caption{Llama3.1-8B, ShareGPT}
\end{subfigure}
\begin{subfigure}[t]{0.22\textwidth}
\includegraphics[width=\linewidth]{figures/e2e/gen_speedup_llama3-8b_instructcoder.pdf}
\caption{Llama3.1-8B, InstructCoder}
\end{subfigure}

\begin{subfigure}[t]{0.22\textwidth}
\includegraphics[width=\linewidth]{figures/e2e/gen_speedup_llama3-70b_gsm8k.pdf}
\caption{Llama3-70B, GSM8K}
\end{subfigure}
\begin{subfigure}[t]{0.22\textwidth}
\includegraphics[width=\linewidth]{figures/e2e/gen_speedup_llama3-70b_cnndailymail.pdf}
\caption{Llama3-70B, CNN/Dailymail}
\end{subfigure}
\begin{subfigure}[t]{0.22\textwidth}
\includegraphics[width=\linewidth]{figures/e2e/gen_speedup_llama3-70b_sharegpt.pdf}
\caption{Llama3-70B, ShareGPT}
\end{subfigure}
\begin{subfigure}[t]{0.22\textwidth}
\includegraphics[width=\linewidth]{figures/e2e/gen_speedup_llama3-70b_instructcoder.pdf}
\caption{Llama3-70B, InstructCoder}
\end{subfigure}

\begin{subfigure}[t]{0.22\textwidth}
\includegraphics[width=\linewidth]{figures/e2e/gen_speedup_qwen-8b-no-thinking_gsm8k.pdf}
\caption{Qwen3-8B, GSM8K}
\end{subfigure}
\begin{subfigure}[t]{0.22\textwidth}
\includegraphics[width=\linewidth]{figures/e2e/gen_speedup_qwen-8b-no-thinking_cnndailymail.pdf}
\caption{Qwen3-8B, CNN/Dailymail}
\end{subfigure}
\begin{subfigure}[t]{0.22\textwidth}
\includegraphics[width=\linewidth]{figures/e2e/gen_speedup_qwen-8b-no-thinking_sharegpt.pdf}
\caption{Qwen3-8B, ShareGPT}
\end{subfigure}
\begin{subfigure}[t]{0.22\textwidth}
\includegraphics[width=\linewidth]{figures/e2e/gen_speedup_qwen-8b-no-thinking_instructcoder.pdf}
\caption{Qwen3-8B, InstructCoder}
\end{subfigure}

\caption{End-to-end performance on non-reasoning workloads. For the $n$-gram method, we report results for both three-token and five-token proposals.}
\label{fig:e2e-speedup}
\end{center}
\end{figure*}

\textbf{End-to-end speedup.} We first present the end-to-end throughput comparison in \autoref{fig:e2e-speedup}.
Across most configurations, SD improves throughput over the baseline without SD. The improvement is most pronounced with small/medium
batch sizes, when the system is memory-bound, so lots of compute can be used for proposing and verification.


\begin{quoting}[leftmargin=1.5em, rightmargin=1.5em, vskip=0pt]
\textbf{Batch Size.}
\emph{Increasing batch size improves absolute throughput but systematically reduces the relative speedup of speculative decoding, an effect that is amplified for larger models.}
\end{quoting}

Consistent with prior work~\cite{liu2024dsd}, larger batch sizes yield higher absolute throughput but lower relative speedup. For example, for the Llama3.1-8B model on the GSM8K dataset, the speedup of EAGLE over the non-SD baseline decreases from 1.73$\times$ to 1.21$\times$ as the batch size increases from 1 to 128.

At a high level, SD trades additional computation for increased throughput or reduced latency. When the batch size is small, the system has sufficient idle compute resources, so the extra computation-spent on proposing and verifying tokens that are ultimately rejected—has limited impact. However, as the batch size grows, the system becomes increasingly compute-bound, making the overhead spent on proposing and verifying rejected tokens more costly, and thus leading to smaller speedups. As shown in \autoref{sec:breakdown}, the verification stage dominates the overall execution time under such conditions.

We further observe that this trade-off between computation and speedup becomes more obvious as the model size increases. For instance, with EAGLE, on the ShareGPT dataset, increasing the batch size from 1 to 32 reduces the speedup by 4.3\% (1.68$\times$ to 1.61$\times$) for Llama3.1-8B, whereas for Llama3-70B, the speedup drops by 14.0\% (1.96$\times$ to 1.72$\times$). This trend arises because the 70B model is executed on four GPUs, so even with small or medium batch sizes, the system is already compute-bound, leaving limited spare compute to verify tokens that are ultimately rejected.


\begin{quoting}[leftmargin=1.5em, rightmargin=1.5em, vskip=0pt]
\textbf{SD Variants and dataset.}
\emph{
The $n$-gram method is generally less effective than other SD approaches across most workloads, with the notable exception of code-editing tasks. Moreover, the draft-model-based method achieves the best performance on the 70B target model; however, the effectiveness diminishes as the target model size decreases (e.g., 8B).}
\end{quoting}
GSM8K, CNN/DailyMail, and ShareGPT all show similar trends, with EAGLE-3 or draft-model-based method achieving the highest speedup across batch sizes. 
In contrast, on InstructCoder, the performance gap among SD variants narrows. For both Llama3.1-8B and Qwen3-8B, the $n$-gram method even outperforms EAGLE and EAGLE-3. This behavior arises because code-editing workloads such as InstructCoder exhibit strong token reuse, which can be effectively exploited by simple drafting methods like $n$-gram.
We further analyze the performance of $n$-gram in \autoref{subsect:ngram-bleu-score}.

% \jerry{NOTE: I updated the data for 70B draft model.}

% Next, we observe that draft-model-based methods achieve performance comparable to EAGLE on the Llama3-70B model, whereas on Qwen3-8B their performance is consistently worse than EAGLE-3 and, in some cases, even worse than $n$-gram. 
Next, we observe that draft-model–based methods outperform EAGLE on the Llama3-70B model in most settings. As shown in \autoref{fig:acc-len-request-level-box-plot}, draft-model–based approaches achieve higher token acceptance rates, while their proposing overhead is comparable to that of EAGLE (see \autoref{fig:exec-breakdown}). This behavior is expected: draft models are pretrained on large-scale corpora, whereas the EAGLE head is only fine-tuned on a relatively limited dataset, which naturally constrains its acceptance rate.

In contrast, on Qwen3-8B, draft-model–based speculative decoding consistently underperforms EAGLE-3 and, in some cases, even the training-free $n$-gram method.
% This degradation is primarily driven by the higher proposing overhead of draft-model–based SD on smaller target models. 
Specifically, for Llama3-70B, we use Llama3.2-1B as the draft model, whereas for Qwen3-8B we use Qwen3-0.6B. Although acceptance rates can be similar across the 70B and 8B setups on the same dataset—or even higher for the 8B case sometimes—draft-model–based SD is nonetheless noticeably less effective on the smaller model. For example, on GSM8K, the acceptance rate is 75\% for the 70B setup and 81\% for the 8B setup, yet the overall performance gain is substantially lower for Qwen3-8B.

This discrepancy arises because the proposing overhead is substantially larger for the 8B model. Specifically, the execution-time ratio between the draft and target model's single forward pass is around 12.5\% for the 70B setup, compared to an estimated 37.5\% for the 8B setup. As a result, drafting becomes significantly more expensive relative to verification on smaller models, which diminishes the overall benefit of draft-model-based SD.

% \jerry{(Hmm, the acceptance rate can be different for different datasets; for example, on GSM8K, it is 55\% for 70B and 80\% for 8B -\>, thus we see draft-model has similar performance on GSM8K for 70B and 8B; For ShareGPT, 53\% for 70B v.s. 60\% for 8B.}, 



% \jerry{We also find that draft-model-based SD maintains similar speedup as EAGLE for Llama3-70B, but its speedup is lower than EAGLE3 for Qwen3-8B.  
% [JERRY-TODO] due to limited choice of draft model, parameter size ratio -> also, drafting overheads; acceptance rates}
% For Llama3-70B on GSM8K, EAGLE’s speedup drops from 1.90$\times$ to 1.00$\times$, while draft-model-based SD declines more gradually from 1.86$\times$ to 1.39$\times$. This can be attributed to its higher average acceptance rate (0.55 versus 0.45 for EAGLE), particularly beneficial at large batch sizes where the verification phase dominates execution time and the system becomes compute-bound.}
% \jerry{@Lily But acceptance rate for draft-model is 1.21x higher but speedup is 1.39x higher. It could be EAGLE's 0.45 acc rate is not very accurate; this is from a quick run over 50 requests.}

% \jerry{We also find that the draft-model-based method maintains similar speedup as EAGLE for Llama3-70B, but its speedup is lower than EAGLE3 for Qwen3-8B. }

% Figures
\begin{figure*}[h!]
\begin{center}
\captionsetup{justification=centering}
\includegraphics[width=0.8\linewidth]{figures/e2e/draft-tree-500reqs/legend.pdf}

\begin{subfigure}[t]{0.22\textwidth}
\includegraphics[width=\linewidth]{figures/e2e/draft-tree-500reqs/q3-8b/sglang-only_gsm8k.pdf}
\caption{Qwen3-8B, GSM8K}
\end{subfigure}
\begin{subfigure}[t]{0.22\textwidth}
\includegraphics[width=\linewidth]{figures/e2e/draft-tree-500reqs/q3-8b/sglang-only_cnndailymail.pdf}
\caption{Qwen3-8B, CNN/Dailymail}
\end{subfigure}
\begin{subfigure}[t]{0.22\textwidth}
\includegraphics[width=\linewidth]{figures/e2e/draft-tree-500reqs/q3-8b/sglang-only_sharegpt.pdf}
\caption{Qwen3-8B, ShareGPT}
\end{subfigure}
\begin{subfigure}[t]{0.22\textwidth}
\includegraphics[width=\linewidth]{figures/e2e/draft-tree-500reqs/q3-8b/sglang-only_instructcoder.pdf}
\caption{Qwen3-8B, InstructCoder}
\end{subfigure}

% L3-70B
\begin{subfigure}[t]{0.22\textwidth}
\includegraphics[width=\linewidth]{figures/e2e/draft-tree-500reqs/l3-70b/sglang-only_gsm8k.pdf}
\caption{Llama3-70B, GSM8K}
\end{subfigure}
\begin{subfigure}[t]{0.22\textwidth}
\includegraphics[width=\linewidth]{figures/e2e/draft-tree-500reqs/l3-70b/sglang-only_cnndailymail.pdf}
\caption{Llama3-70B, CNN/Dailymail}
\end{subfigure}
\begin{subfigure}[t]{0.22\textwidth}
\includegraphics[width=\linewidth]{figures/e2e/draft-tree-500reqs/l3-70b/sglang-only_sharegpt.pdf}
\caption{Llama3-70B, ShareGPT}
\end{subfigure}
\begin{subfigure}[t]{0.22\textwidth}
\includegraphics[width=\linewidth]{figures/e2e/draft-tree-500reqs/l3-70b/sglang-only_instructcoder.pdf}
\caption{Llama3-70B, InstructCoder}
\end{subfigure}
\caption{End-to-end performance on non-reasoning workloads for tree-style verification. We compare a chain setting ($k{=}3$) against tree settings with $k{=}6$ and $k{=}21$, using a fixed depth of 3 for all runs.}
\label{fig:e2e-speedup-tree}
\end{center}
\end{figure*}

\begin{quoting}[leftmargin=1.5em, rightmargin=1.5em, vskip=0pt]
\textbf{Tree-Style Verification.}
\emph{Tree-based methods provide slightly higher speedup at batch size 1, but the advantage quickly disappears as batch size increases. Across batch sizes, the chain-style method often remains the more performant.}
\end{quoting}

% We study tree-style verification in SGLang v0.5.9~\cite{zheng2024sglang} using Qwen3-8B and Llama3-70B on the non-reasoning workloads. We use SGLang here because the current vLLM draft-tree path we tested is not yet fully optimized, making it difficult to compare tree settings fairly. To keep the comparison controlled, we fix the tree depth to 3 in all runs. We use the chain setting with $k{=}3$ as the baseline, evaluate $k{=}21$ with a branching factor of 4 to roughly match the tree structure used in the EAGLE paper~\cite{li2024eagle}, and include $k{=}6$ with a branching factor of 2 as a smaller intermediate tree. Here, $k$ denotes the number of draft tokens verified in parallel by the target model. Unless otherwise noted, we use the same experimental settings as described earlier. By default, SGLang enables a dynamic draft-tree policy that adjusts the tree structure during decoding~\cite{li2024eagle}. Also, FlashAttention3~\cite{shah2024flashattention} is used for the chain case, while FlashInfer~\cite{ye2025flashinfer} is used for the tree cases by default.

We next evaluate tree-style verification in SGLang v0.5.9~\cite{zheng2024sglang} using Qwen3-8B and Llama3-70B on the non-reasoning workloads. We use SGLang for this study because the draft-tree path in the vLLM version we evaluated is not yet sufficiently optimized for a fair comparison across tree configurations. To control the comparison, we fix the tree depth to 3 in all runs. We use chain-style verification with $k{=}3$ as the baseline, evaluate a wider tree with $k{=}21$ and branching factor 4 to approximately match the tree structure used in EAGLE~\cite{li2024eagle}, and include an intermediate tree with $k{=}6$ and branching factor 2. Here, $k$ denotes the number of draft tokens verified in parallel by the target model. Unless otherwise noted, all other settings follow those described earlier. By default, SGLang uses a dynamic draft-tree policy which optimises the tree structure during decoding~\cite{li2024eagle}. In this setup, FlashAttention-3~\cite{shah2024flashattention} is used for the chain configuration, whereas FlashInfer~\cite{ye2025flashinfer} is used for the tree configurations.

As shown in \autoref{fig:e2e-speedup-tree}, tree-based EAGLE/EAGLE-3 achieves slightly higher speedup than the chain-based settings at batch size 1. For example, on Qwen3-8B with GSM8K, speedup increases from 1.65$\times$ for the chain to 1.68$\times$ for $k{=}6$ and 1.85$\times$ for $k{=}21$. On Llama3-70B with ShareGPT, it rises from 1.81$\times$ to 1.90$\times$ and 2.03$\times$. However, this benefit quickly disappears as batch size increases. By batch size 64, the $k{=}21$ tree falls below 1$\times$ speedup on all workloads for both models, whereas the chain remains above 1$\times$ throughout.

This is likely because wider trees increase the accepted length, but also verify many more tokens that are later rejected. On Qwen3-8B with GSM8K, the accepted length increases from 2.25 for the chain to 2.51 and 2.92 for the $k{=}6$ and $k{=}21$ trees, respectively, while the acceptance rate decreases from 0.415 to 0.300 and 0.095. We observe the same trend on Llama3-70B with ShareGPT (accepted length: 2.29 $\rightarrow$ 2.55 $\rightarrow$ 2.93; acceptance rate: 0.429 $\rightarrow$ 0.310 $\rightarrow$ 0.097). As discussed in \autoref{sec:breakdown}, verification dominates the execution cost, so the overhead of these additional rejected branches becomes increasingly expensive at larger batch sizes. Thus, tree-style verification provides only a narrow low-batch-size benefit, while chain-style verification remains the more robust setting overall.

\begin{figure*}[h!]
\captionsetup{justification=centering}
\begin{center}
\begin{subfigure}[t]{0.22\textwidth}
\includegraphics[width=\linewidth]{figures/e2e/thinking_qwen-8b-thinking_gsm8k.pdf}
\caption{Qwen3-8B-Thinking, GSM8K}
\end{subfigure}
\begin{subfigure}[t]{0.22\textwidth}
\includegraphics[width=\linewidth]{figures/e2e/thinking_qwen-8b-thinking_cnndailymail.pdf}
\caption{Qwen3-8B-Thinking, CNN/Dailymail}
\end{subfigure}
\begin{subfigure}[t]{0.22\textwidth}
\includegraphics[width=\linewidth]{figures/e2e/thinking_qwen-8b-thinking_sharegpt.pdf}
\caption{Qwen3-8B-Thinking, ShareGPT}
\end{subfigure}
\begin{subfigure}[t]{0.22\textwidth}
\includegraphics[width=\linewidth]{figures/e2e/thinking_qwen-8b-thinking_instructcoder.pdf}
\caption{Qwen3-8B-Thinking, InstructCoder}
\end{subfigure}

\begin{subfigure}[t]{0.22\textwidth}
\includegraphics[width=\linewidth]{figures/e2e/thinking_qwen-8b-thinking_aime.pdf}
\caption{Qwen3-8B-Thinking, AIME22-24}
\end{subfigure}
\begin{subfigure}[t]{0.22\textwidth}
\includegraphics[width=\linewidth]{figures/e2e/thinking_qwen-8b-thinking-gpqa_gpqa_main.pdf}
\caption{Qwen3-8B-Thinking, GPQA-Main}
\end{subfigure}
\begin{subfigure}[t]{0.22\textwidth}
\includegraphics[width=\linewidth]{figures/e2e/thinking_glm4.5_aime.pdf}
\caption{GLM-4.5-Air, AIME22-24}
\end{subfigure}
\begin{subfigure}[t]{0.22\textwidth}
\includegraphics[width=\linewidth]{figures/e2e/thinking_glm4.5_gpqa_main.pdf}
\caption{GLM-4.5-Air, GPQA-Main}
\end{subfigure}
\caption{End-to-end performance on reasoning workloads. }
\label{fig:e2e-reasoning}
\end{center}
\end{figure*}

\begin{quoting}[leftmargin=1.5em, rightmargin=1.5em, vskip=0pt]
\textbf{Reasoning Workloads.}
\emph{Models and drafting methods that sustain high acceptance rates over long contexts—such as EAGLE-3 and, in some cases, $n$-gram—derive the greatest benefit on reasoning tasks. In contrast, MTP performance is limited by the reuse of a single MTP head across drafted tokens.}
\end{quoting}

As reasoning workloads become increasingly prevalent, understanding the effectiveness of speculative decoding (SD) in this regime becomes important. These workloads feature short prompts but long generation sequences. We evaluate SD on two representative reasoning benchmarks, AIME22–24~\cite{AIME-HF} and GPQA-Main~\cite{rein2023gpqagraduatelevelgoogleproofqa}.

As shown in \autoref{fig:e2e-reasoning}, we report the absolute generation throughput across different batch sizes. To ensure fair and consistent comparisons, we restrict our evaluation to medium batch sizes. Reasoning tasks typically involve long generation sequences, which can trigger request preemption once the KV cache becomes full. By avoiding such preemption, the reported results more accurately reflect the inherent efficiency of speculative decoding rather than performance artifacts caused by memory pressure.

For Qwen3-8B-Thinking, EAGLE-3 consistently leads across datasets, achieving 1.64$\times$--1.80$\times$ speedup over on GPQA-Main and AIME22–24, whereas $n$-gram performs comparably at 1.50$\times$--1.58$\times$. This trend likely arises because the model produces longer contexts and repetitive symbolic patterns, increasing opportunities for reuse. As shown later in~\autoref{sec:acc-behavior}, the average acceptance length in GPQA-Main generally grows with sequence length, with $n$-gram’s acceptance rising faster than EAGLE-3’s.

For GLM-4.5-Air, MTP outperforms $n$-gram but falls short of the theoretical upper bound, achieving 1.3$\times$--1.8$\times$ speedup on GPQA-Main. This gap is likely because the released open-source weights include only the first MTP module, which is trained primarily to predict the first token accurately but reused autoregressively for subsequent ones. Consequently, its position-wise acceptance rate (\(0.92 \rightarrow 0.68 \rightarrow 0.38\) on GPQA-Main), accuracy declines sharply across drafted tokens, shrinking the acceptance rate and thereby constraining the achievable speedup.



%% <<< end performance.tex


%% >>> begin understanding_perf.tex
\section{Understanding the Performance}
\label{sec:understand-perf}
In this section, we aim to have a better understanding of the end-to-end speedups reported in \autoref{sect:e2e-perf}. 
We start by breaking down where time and memory are spent on different SD stages -- drafting, verification, rejection sampling, and others, and how these keep us from reaching the theoretical upper bound. We then examine how acceptance patterns within a request, across requests, and across datasets. 

% While the $n$-gram approach, as a training-free SD variant, typically underperforms compared with learned SD methods in \autoref{sect:e2e-perf}, its acceptance patterns indicate untapped potential; in particular, we observe surprisingly long acceptance segments under specific conditions. This motivates a case study to investigate when this lightweight, training-free method can be especially advantageous.
% , linking prompt–output overlap, measured by BLEU-$n$ score~\cite{papineni2002bleu}, to the strong speedups achieved by the $n$-gram method. 

\subsection{End to End Speedup Formula}
As pointed in the original SD paper~\cite{leviathan2023fast},
Let $c$ denote the execution-time ratio of a single forward pass between the drafting method and the target model, $\alpha$ is the token acceptance rate and $k$ is the proposed length. The expected overall wall-time speedup achieved by speculative decoding, can be expressed as follows:\footnote{This formulation assumes that the $k + 1$ simultaneous evaluations of the target model take approximately the same amount of time as generating a single token in parallel. This assumption typically holds when the system is memory-bound, which is common at small batch sizes, but may break down once the system becomes compute-bound.}
\begin{equation}
\label{eq:speedup}
     E(speedup)=\frac{1-\alpha^{k+1}}{(1-\alpha)(kc+1)}.
\end{equation}
As shown, the speedup depends on two factors: the proposed length ($k$), the execution-time ratio between the drafting method and the target model ($c$), and the token acceptance rate ($\alpha$). Importantly, the absolute execution time of the drafting method is irrelevant; only its runtime relative to the target model affects the achievable speedup.

Employing a more sophisticated drafting method can improve the token acceptance rate $\alpha$, but often at the cost of increased execution time, leading to a larger $c$. Consequently, optimal speedup is achieved by carefully choosing a drafting method that balances the trade-off between $c$ and $\alpha$. In the following sections, we evaluate these factors and additionally analyze the memory overhead of SD.

\section{Execution Time and Memory Breakdown}
\label{sec:breakdown}
% \jerry{
% 1. describe figures
% 2. does percentage vary across model sizes (different TPs)?
% 3. verification cost is the majority
% }
% \lily{
% Old comments, migrated:
% 1. explain the speedup difference between measured ones in vLLM and eagle repo.
% 2. what's optimal speedup if we get rid of engineering overhead. explain we already cudagraph the propose/verify model, it's efficient.
% 3. add results of 70B model, add communication time.
% 4. discuss the potential of async scheduling.}

% \jerry{Should we explain why we didn't analyse time breakdown for draft-model based and MTP?}
\subsection{Execution Time}
% \lily{(1)We should include a detailed execution-time breakdown for draft-model–based methods in Section4.} 
In \autoref{fig:exec-breakdown}, we present the execution time breakdown of different stages in the baseline LLM execution and SD settings. For all experiments, we sample 500 requests from CNN/DailyMail. All other runtime settings are the same as \autoref{tab:exp-config}. 
For each configuration, we report the fraction of total execution time attributable to four components:
\textit{(i)}~\textbf{Drafting}, the time to generate speculated tokens, which corresponds to the lookup operation for $n$-gram and the autoregressive generation of the EAGLE heads for EAGLE and EAGLE-3;
\textit{(ii)}~\textbf{Verification}, the time spent by the target model to validate proposed tokens;
\textit{(iii)}~\textbf{Rejection Sampling}, time spent on generating final tokens based on verification logits;
and \textit{(iv)}~\textbf{Other Overheads}, all other overheads in the vLLM execution system. 

\begin{figure*}
\centering
\begin{subfigure}[t]{0.43\linewidth}
\includegraphics[width=\linewidth]{figures/time_breakdown/legend_only.pdf}
\end{subfigure}

\begin{subfigure}[t]{0.33\linewidth}
\includegraphics[width=\linewidth]{figures/time_breakdown/cnndailymail_llama3.1-8B_time_breakdown_legend_False.pdf}
\caption{Llama3.1-8B}
\end{subfigure}
\begin{subfigure}[t]{0.33\linewidth}
\includegraphics[width=\linewidth]{figures/time_breakdown/cnndailymail_Llama-3-70B_time_breakdown_legend_False.pdf}
\caption{Llama3-70B}
\end{subfigure}
\begin{subfigure}[t]{0.33\linewidth}
\includegraphics[width=\linewidth]{figures/time_breakdown/cnndailymail_Qwen-3B_time_breakdown_legend_False.pdf}
\caption{Qwen3-8B}
\end{subfigure}
\caption{Execution time breakdown across different models.}
\label{fig:exec-breakdown}
\vspace{-0.1in}
\end{figure*}


\begin{quoting}[leftmargin=1.5em, rightmargin=1.5em, vskip=0pt]
\textbf{Execution Time.}
% \emph{Verification time dominates across all speculative decoding variants, while sampling time is negligible. Drafting time is negligible for the $n$-gram method and remains below 20\% for EAGLE.}
\emph{Verification time generally accounts for the largest share of runtime across speculative decoding variants, while sampling time is negligible. Drafting is negligible for $n$-gram, stays below 20\% for EAGLE/EAGLE3, and stays under 40\% for draft-model methods acorss all batch sizes.}
\end{quoting}

We observe that the verification stage generally takes the largest execution time, ranging from 42\% to 95\% in all execution methods, which increases with the model size and batch size. 

Drafting costs, however, vary substantially by proposal mechanism. For $n$-gram, drafting time accounts for less than 2\% across all executions, adding negligible overhead to the baseline execution. Both EAGLE and EAGLE3 show similar drafting overhead characteristics, accounting for a 12\% to 20\% of the execution time at batch size 1, and 3\% to 7\% at batch size of 512.
Draft-model-based methods differ in that they must run a separate smaller model autoregressively to propose tokens, leading to a larger drafting overhead when the batch size is small. For Llama3-70B, drafting takes from 21\% at batch size of 1 to 3\% at batch size of 512. For Qwen3-8B, drafting is about 47\% at batch size of 1 and decreases to 16\% at batch size of 512. As a result, the verification fraction is correspondingly lower (roughly 72\% for Llama3-70B and 42\% for Qwen3-8B at batch size of 1). 
% Verification still remains the dominant component in general, except for Qwen3-8B at batch size of 1, where drafting is slightly larger.

% For drafting, $n$-gram drafting time accounts for less than 2\% across all executions, adding negligible overhead to the baseline execution. Both EAGLE and EAGLE3 show similar drafting overhead characteristics, accounting for a 12\% to 20\% of the execution time at batch size 1, and 3\% to 7\% at batch size of 512.
% As such, even if EAGLE-based speculative decoding may provide a better acceptance rate, it may be beneficial to use $n$-gram over EAGLE when the model or the batch size is small, as shown in more detail later.
Furthermore, the sampling stage takes a negligible amount of time, accounting for less than 1.7\% of total time across all executions. 
% While the vLLM's system overhead increases slightly with the batch size, the increase is marginal compared to the increase in the verification time. 
Lastly, vLLM overhead accounts for 3–12\% of total execution time. This fraction decreases as model and batch size grow, since fixed overheads—such as scheduling—are independent of model scale and therefore become amortized at larger batch size or model size.

\textbf{Implication.} Verification cost dominates the end-to-end execution of speculative decoding, indicating that running the large target model remains the primary computational bottleneck. When the proposed tokens are ultimately accepted, this cost is well justified. However, verifying tokens that are later rejected can incur substantial wasted computation. This observation motivates a natural question: how much of the verification compute is truly useful? Conversely, what speedup could be achieved if verification were performed only on “correct” tokens, eliminating wasted work? These questions form the basis of \autoref{sec:theoretical-speedup}, where we estimate the theoretical upper bound for SD.

\subsection{Memory}
In \autoref{tab:memory}, we present the GPU memory usage of each SD variant with Llama3.1-8B, Llama3-70B and Qwen3-8B.
% , excluding the vLLM system overhead, such as memory fragmentation in KV cache blocks. 
We use FP-16 precision and report the static memory overhead from the model weights and per-token KV cache calculated using model specs. We do not report intermediate memory overheads, such as the activation tensors during execution. The detailed calculation steps can be found in \autoref{appendix:mem}.

\begin{table}[t] % Use [htbp] for better float placement
    \small
    % --- First Sub-table ---
    \begin{subtable}[t]{\columnwidth}
        \centering
        \begin{tabular}{c|c|c|c}
             \toprule
              & No SD/$n$-gram & EAGLE & EAGLE 3\\
             \midrule        
             Static (GiB) & 14.96 & 15.43 & 15.75 \\
             Per-token (KiB) & 128 & 132 & 132 \\
             \bottomrule
        \end{tabular}
        \caption{Memory breakdown on Llama3.1-8B} % Optional caption for sub-table (a)
        \label{tab:sub1} % Label for the first sub-table
    \end{subtable}
    \hfill % This command pushes the two subtables apart
    % --- Second Sub-table ---
    \begin{subtable}[t]{\columnwidth}
        \centering
        \begin{tabular}{c|c|c|c}
              \toprule
              & No SD/$n$-gram & EAGLE & D-M (1B)\\
             \midrule        
             Static (GiB) & 131.4 & 133.3 & 133.7 \\
             Per-token (KiB) & 320 & 324 & 352\\
             \bottomrule
        \end{tabular}
        \caption{Memory breakdown on Llama3-70B} 
        \label{tab:sub2} % Label for the second sub-table
    \end{subtable}
    \hfill
    \begin{subtable}[t]{\columnwidth}
        \centering
        \begin{tabular}{c|c|c|c}
              \toprule
              & No SD/$n$-gram & EAGLE3 & D-M (0.6B)\\
             \midrule        
             Static (GiB) & 15.25 & 16.00 & 16.37 \\ 
             Per-token (KiB) & 144 & 148 & 256\\
             \bottomrule
        \end{tabular}
        \caption{Memory breakdown on Qwen3-8B} 
        \label{tab:sub3}
    \end{subtable}
    
    \caption{Static and per-token memory usage of the baseline model execution and different speculative decoding methods. DM stands for Draft-Model-based method.}
    \label{tab:memory} % The main label for the combined table
    \vspace{-0.1in}
\end{table}


\begin{quoting}[leftmargin=1.5em, rightmargin=1.5em, vskip=0pt]
\textbf{Memory.}
\emph{Speculative decoding generally incurs minimal memory overhead, including both static parameters and KV cache: it is zero for $n$-gram, below 10\% for EAGLE/EAGLE-3, while draft-model-based methods incur overhead that depends on the chosen draft model.}
\end{quoting}
% We also do not report memory overhead required for $n$-gram token lookup, as it is done entirely on the CPU.
We find that the memory overhead of speculative decoding (SD) is minimal. For $n$-gram SD, draft tokens are sampled from CPU-resident generation history, incurring no GPU memory overhead. 

\textbf{Static Memory.} Static memory corresponds to the memory required to store model weights. Certain SD variants add extra layers or even an independent model, thereby incurring additional static memory overhead. EAGLE-based SD introduces an additional Transformer layer and thus requires extra static weights, but the overhead is small: 3.1\% (EAGLE)/ 5.3\% (EAGLE3) for Llama-3.1-8B, 1.4\% (EAGLE) for Llama-3-70B, and 4.9\% (EAGLE3) for Qwen3-8B.
For draft-model-based SD, static memory overhead depends on the size of the draft model; in our setup, pairing Llama-3-70B with Llama-3.2-1B incurs a 1.8\% increase while pairing Qwen3-8B with Qwen3-0.6B incurs a more significant 7.3\% increase. 

\textbf{Per-token Memory.} Per-token memory refers to the size of the KV cache allocated for each generated token. The total KV cache footprint grows with the number of generated but unfinished requests. Consequently, for workloads with long generation lengths such as reasoning workloads, memory cost introduced by KV cache becomes a dominant factor in overall memory consumption.
Per-token overhead is modest for EAGLE-based methods, as each EAGLE layer requires the same KV cache as a single attention layer (4 KiB per token), resulting in a 3.1\% overhead for Llama-3.1-8B and 1.3\% for Llama-3-70B. 
In contrast, draft-model–based SD incurs substantially higher per-token memory overhead, as it relies on an additional multi-layer model for token proposal. In our configuration, pairing a 0.6B draft model with an 8B target model increases the per-token memory footprint from 144 KiB to 256 KiB, corresponding to a 1.77$\times$ increase.



% \textbf{Observation.} \lily{TODO, small bach size, enough memory.}



% \begin{figure*}[h]
%     \centering
%     \includegraphics[width=0.24\linewidth]{figures/time_breakdown/cnndailymail_llama3.1-8b_none_128_time_ratios.pdf}
%     \includegraphics[width=0.24\linewidth]{figures/time_breakdown/cnndailymail_llama3.1-8b_ngram_128_time_ratios.pdf}
%     \includegraphics[width=0.24\linewidth]{figures/time_breakdown/cnndailymail_llama3.1-8b_eagle_128_time_ratios.pdf}
%     \includegraphics[width=0.24\linewidth]{figures/time_breakdown/cnndailymail_llama3.1-8b_eagle3_128_time_ratios.pdf}
%     \caption{Histogram of execution time breakdown for Llama3.1-8B (TP=1) model.}
    
%     \includegraphics[width=0.3\linewidth]{figures/time_breakdown/cnndailymail_Qwen_Qwen3-8B_none_128_time_ratios.pdf}
%     \includegraphics[width=0.3\linewidth]{figures/time_breakdown/cnndailymail_Qwen_Qwen3-8B_ngram_128_time_ratios.pdf}
%     \includegraphics[width=0.3\linewidth]{figures/time_breakdown/cnndailymail_Qwen_Qwen3-8B_eagle_128_time_ratios.pdf}
%     \caption{Histogram of execution time breakdown for Qwen3-8B (TP=1) model.}

%     \includegraphics[width=0.3\linewidth]{figures/time_breakdown/cnndailymail_meta-llama_Meta-Llama-3-70B-Instruct_none_128_time_ratios.pdf}
%     \includegraphics[width=0.3\linewidth]{figures/time_breakdown/cnndailymail_meta-llama_Meta-Llama-3-70B-Instruct_ngram_128_time_ratios.pdf}
%     \includegraphics[width=0.3\linewidth]{figures/time_breakdown/cnndailymail_meta-llama_Meta-Llama-3-70B-Instruct_eagle_128_time_ratios.pdf}
%     \caption{Histogram of execution time breakdown for Llama3-70B (TP=4) model.Time ratio: time spent on the given phase divided by the total execution time of the step. \jerry{change histogram to a bar chart!}}

%     \label{fig:exec-breakdown}
% \end{figure*}

%% <<< end understanding_perf.tex


%% >>> begin acceptance_behavior.tex
\begin{figure}[h!]
  \centering
  \begin{subfigure}{0.45\linewidth}
    \includegraphics[width=\linewidth]{figures/acc_len/heatmap/instructcoder_ngram_llama3.1-8B_acceptance_stats.pdf}
    \caption{InstructCoder, $n$-gram}
  \end{subfigure}
  \hspace{1pt}
  \begin{subfigure}{0.45\linewidth}
    \includegraphics[width=\linewidth]{figures/acc_len/heatmap/sharegpt_ngram_llama3.1-8B_acceptance_stats.pdf}
    \caption{ShareGPT, $n$-gram}
  \end{subfigure}

   \begin{subfigure}{0.45\linewidth}
    \includegraphics[width=\linewidth]{figures/acc_len/heatmap/instructcoder_eagle_llama3.1-8B_acceptance_stats.pdf}
    \caption{InstructCoder, EAGLE}
  \end{subfigure}
  \hspace{1pt}
  \begin{subfigure}{0.45\linewidth}
    \includegraphics[width=\linewidth]{figures/acc_len/heatmap/sharegpt_eagle_llama3.1-8B_acceptance_stats.pdf}
    \caption{ShareGPT, EAGLE}
  \end{subfigure}

  \caption{Generation length per token position for Llama3.1-8B. Requests are sorted based on generation length. Darker colors indicate that more tokens are generated at the corresponding position.}
  \label{fig:acceptance-length-heatmap}
\end{figure}



\section{Acceptance Behavior}
\label{sec:acc-behavior}
To understand the acceptance rate, we measure the number of draft tokens generated\footnote{Number of generated tokens = accepted tokens proposed by the drafting method + 1 bonus token. Similarly, for generation length, it is computed by adding the acceptance length by one.} by the target model at each decoding step. 
For each dataset, we sample up to 200 requests. Each SD method proposes up to 20 tokens per step, and we record how many of these are generated at that step.  
After each generation step, even if multiple tokens were generated, we discard all but one and advance to the next position. 
Essentially, we approximate the maximum number of tokens that could be accepted at each generation step.
We analyze results for two non-reasoning models (Llama-3-70B and Llama-3.1-8B; max 512 output tokens) and one reasoning model (Qwen-3-8B-Thinking; max 32K output tokens).


\autoref{fig:acceptance-length-heatmap} visualizes the number of tokens generated at each output position for InstructCoder and ShareGPT (notice SD generates at least one token at each position, even if all proposed tokens are rejected). Each row represents a sampled request and each column a generation position, with color indicating the number of tokens produced by SD. 

\begin{quoting}[leftmargin=1.5em, rightmargin=1.5em, vskip=0pt]
\textbf{Summary.}
\emph{We observe the variance of generated token lengths along three dimensions: within a single request, across requests, and across datasets.}
\end{quoting}
\autoref{fig:acceptance-length-heatmap} gives a high level feeling of the generation length at each position across requests. $n$-grams and EAGLE show different patterns. $n$-grams show high variance, lots of very dark and very light positions, some positions generate a lot, while others only generate one. EAGLE show good average generation length (averge color is darker) and also less variance. But there are variance across all dimentions, which we will describe in more detail below. 


% Acceptance varies notably across positions (horizontal), requests (vertical), and datasets: within a request, acceptance fluctuates along the sequence; across requests, some are consistently easier to speculate than others; across datasets, for example, \autoref{fig:acceptance-length-heatmap} shows InstructCoder generally achieves higher acceptance (darker) than ShareGPT.

% \autoref{fig:acceptance-length-heatmap} presents the number of generated tokens (notice SD generates at least one token at each position, even if all proposed tokens are rejected) at each generation position for InstructCoder and ShareGPT. Each row corresponds to one sampled request, and each column to a generation position, with color indicating the number of tokens generated by SD\footnote{Number of generated tokens = number of accepted tokens proposed by the drafting method + 1}.
% As shown, acceptance behavior exhibits \textbf{significant variance} across positions (horizontal), requests (vertical), and datasets. Within a single request, acceptance fluctuates along the sequence, showing that some positions are easier to predict than others. Across requests, some are consistently easier to speculate (darker rows), while others remain difficult. Across datasets, we observe that, for example, InstructCoder exhibits generally higher acceptance (darker in general) than ShareGPT. These patterns motivate a more structured analysis along three dimensions: acceptance behavior within a request, across requests, and across datasets.


\begin{quoting}[leftmargin=1.5em, rightmargin=1.5em, vskip=0pt]
\textbf{Within A Request.}
\emph{Longer generation workloads yield more accepted tokens per position, with $n$-gram benefiting from repetitive reasoning patterns but losing effectiveness near the end as the generation shifts toward conclusions.}
% This effect is more pronounced for the $n$-gram method.} 

\end{quoting}

Since different requests can have varying generation lengths, we sample 50 requests from each generation-length category: requests with generation length $<$4K, between 4–8K, between 8–13K, and over 13K. 
% \jerry{Note that some requests in the last category reach the 32K output length limit before they can finish, and thus are excluded from the analysis.}
% To ensure fair comparison across outputs of different lengths, we first sort all requests by their output length and divide them into four equal-count bins (50 requests each). 
In \autoref{fig:acc-len-token-level}, token positions are partitioned into ten intervals, and the mean generation length is computed and plotted for each interval.

For reasoning workloads, generation length generally increases as generation progresses for both methods, but the growth is substantially steeper for $n$-gram. As shown in \autoref{fig:acc-len-token-level}, on GPQA-Main with Qwen-3-8B-Thinking, $n$-gram rises from roughly $1.6$--$3.7$ generated tokens across the sequence in short outputs ($<$4K tokens) to $2.7$--$5$ tokens in the longest responses ($>$13K tokens). 
EAGLE-3 grows much more gradually from around $2.1$--$2.5$ among the shortest requests and to roughly $2.5$--$5.7$ tokens among the longest requests. 
These results indicate that long scientific and technical generations accumulate locally repetitive structures, such as recurring variables, equations, and phrases, that $n$-gram can increasingly exploit. On the other hand, \textit{EAGLE}-based methods remain less sensitive to such surface-level recurrence.

For $n$-gram, we additionally observe a decrease towards the last few token position intervals. This is likely because the model shifts from step-by-step reasoning, where phrases and patterns tend to repeat a lot, to writing the final answer and stopping, which is less repetitive and therefore harder for $n$-gram to reuse previous context. Thus, in the last few token intervals, more requests are in the final-answer stage, so the average generation length drops.




% \jerry{Notably, requests with generation length over 13K do not show this drop. This is because 24 out of 50 of these generations hit the 32K limit before they finish, so the last token intervals still contain ongoing reasoning rather than a finalized answer and stopping behavior.}

% \jerry{@Lily, I feel the explanation for 13K sounds more like a flaw in the choice of the requests, rather than a finding or insight, I should filter out these not completed requests. -$>$ I have verified that filtering out these uncompleted requests will make ngram in $>$13k category show the same increase-then-decrease pattern. --$>$ Thus, I have updated the figures and plan to remove this part.}

\begin{figure} [h!]
\captionsetup{justification=centering}
  \centering
  \begin{subfigure}{0.45\linewidth}
    \includegraphics[width=\linewidth]{figures/acc_len/token_level/overlaid_binned_positions_gpqa_main_Qwen3-8B_bin0_range0-4190.pdf}
    \caption{Requests with outputs shorter than 4K tokens}
  \end{subfigure}
    \hspace{1pt}
    \begin{subfigure}{0.45\linewidth}
    \includegraphics[width=\linewidth]{figures/acc_len/token_level/overlaid_binned_positions_gpqa_main_Qwen3-8B_bin3_range13351-32771.pdf}
    \caption{Requests with outputs longer than 13K tokens}
  \end{subfigure}
  \begin{subfigure}{0.45\linewidth}
    \includegraphics[width=\linewidth]{figures/acc_len/token_level/overlaid_binned_positions_gpqa_main_Qwen3-8B_bin1_range4190-7811.pdf}
    \caption{Requests with outputs of 4K-8K tokens}
  \end{subfigure}
    \hspace{1pt}
    \begin{subfigure}{0.45\linewidth}
    \includegraphics[width=\linewidth]{figures/acc_len/token_level/overlaid_binned_positions_gpqa_main_Qwen3-8B_bin2_range7811-13351.pdf}
    \caption{Requests with outputs of 8K-13K tokens}
  \end{subfigure}
  \vspace{-0.1in}
\caption{Average generation length (tokens) per output token position for Qwen3-8B-Thinking on GPQA-Main. Shaded bands show the 95\% confidence interval.}
  \label{fig:acc-len-token-level}
\end{figure}


\begin{figure} [t!]
  \centering
  \includegraphics[width=0.8\linewidth]{figures/acc_len/request_level/legend.pdf}

  \begin{subfigure}{0.45\linewidth}
    \includegraphics[width=\linewidth]{figures/acc_len/request_level/box_by_model_gsm8k_acc_len.pdf}
    \caption{GSM8K}
  \end{subfigure}
      \hspace{1pt}
    \begin{subfigure}{0.45\linewidth}
    \includegraphics[width=\linewidth]{figures/acc_len/request_level/box_by_model_cnndailymail_acc_len.pdf}
    \caption{CNNDailyMail}
  \end{subfigure}
      \hspace{1pt}
  \begin{subfigure}{0.45\linewidth}
    \includegraphics[width=\linewidth]{figures/acc_len/request_level/box_by_model_sharegpt_acc_len.pdf}
    \caption{ShareGPT}
  \end{subfigure}
      \hspace{1pt}
  \begin{subfigure}{0.45\linewidth}
    \includegraphics[width=\linewidth]{figures/acc_len/request_level/box_by_model_instructcoder_acc_len.pdf}
    \caption{InstructCoder}
  \end{subfigure}
\caption{Request-level generation length across datasets and models. Each box shows the distribution of per-request mean generation length, with the box spanning the 25th–75th percentiles and whiskers covering the 5th–95th percentiles. Models from left to right are Llama3.1-8B, Llama3-70B and Qwen3-8B-Thinking. Results for reasoning datasets (AIME22–24 and GPQA-Main) are shown in~\autoref{fig:acc-len-request-level-box-plot-reasoning}.}
  \label{fig:acc-len-request-level-box-plot}
\end{figure}



\begin{quoting}[leftmargin=1.5em, rightmargin=1.5em, vskip=0pt]
\textbf{Across Requests and Datasets.}
\emph{Draft-model-based method typically achieves the longest median generation lengths across all workloads, EAGLE attains more stable per-request generation lengths, whereas $n$-gram shows higher variance and a heavy-tailed generation-length distribution.}
\end{quoting}

% \jerry{Updated, please read the entire section again.} 
We observe substantial variance in per-request generation length in \autoref{fig:acc-len-request-level-box-plot}. Even within the same workload, the efficiency of speculative decoding can differ drastically. For example, on InstructCoder with Llama3-70B, EAGLE produces relatively short request-level generation lengths (2.7-7.4 tokens), whereas $n$-gram and draft-model-based SD achieve longer lengths with wider spreads (1.1–15.0 tokens and 5.6-18.3 tokens, respectively). In other workloads (e.g. ShareGPT and CNN/DailyMail) and target models (e.g. Qwen3-8B), draft-model-based methods similarly shows higher generation lengths than $n$-gram and EAGLE.

% Too many data, rewrite this.
% We observe that per-request generation length varies widely in \autoref{fig:acc-len-request-level-box-plot}. 
% Even within the same workload, the efficiency of speculative decoding can differ drastically. On InstructCoder with Llama3-70B, request-level generation lengths span from 1.1–15.0 tokens for $n$-gram (mean 7.0, std 4.3), 2.7-7.4 for EAGLE (mean 4.9, std 1.4) and 5.6-18.3 for Draft-Model-based (mean 12.4, std 3.8). Similar disparities appear in ShareGPT (1.0–4.7 vs.\ 2.0–4.2 vs.\ 3.4-10.8, std 1.6 vs.\ 1.0 vs.\ 2.5), and CNN/DailyMail (1.2–4.0 vs.\ 2.1–3.0 vs.\ 3.6-7.0, std 0.9 vs.\ 0.3 vs.\ 1.0).
% and GSM8K (1.0-2.5 vs.\ 2.0–2.7 vs.\ 4.3-12.9, std  0.5 vs.\ 0.2 vs.\ 2.7), 
% where the average generation length of most requests clusters around 1.5–3 tokens for $n$-gram and eagle. 

These distributions reveal distinct proposal behaviors for each SD variant. Draft-model-based methods often achieve substantially longest generation lengths, but their whiskers (5th-95th percentiles) are also broadly spread. 
% \jerry{@lily, does this explanation sound reasonable to you?} 
The reason behind this may be request-dependent alignment: for some requests, the draft model can track the target model closely so that many proposed tokens are accepted, while for others, it diverges early and only a few draft tokens are accepted after verification.

For $n$-gram, it exhibits heavy tails and large variance across requests, with standard deviations roughly 2x-5x higher than those of EAGLE-based methods. 
While most $n$-gram requests produce short accepted spans, a small subset yield exceptionally long bursts (often exceeding 15 tokens), as seen as long blue strips in \autoref{fig:acceptance-length-heatmap}. 
In code-editing tasks, these typically correspond to locally repetitive content (e.g., reused identifiers, function templates, or class definitions). We select and show an example in \autoref{fig:ngram_prompt_output_full}. Such bursty acceptance behavior reveals that $n$-gram speculation relies on discrete pattern matches that may occur infrequently but can produce large payoffs when present. Consequently, $n$-gram's performance is high when the local contexts are repetitive, but it remains unstable across requests and workloads.


In contrast, EAGLE and EAGLE-3 show compact, symmetric distributions centered near their medians (typically 2–4 tokens), indicating steadier acceptance driven by learned contextual representations rather than surface-level overlap.

These distributional behaviors are consistent with the end-to-end outcomes in \autoref{fig:e2e-speedup}. Draft-model-based methods deliver the largest speedups when the cost of running a separate draft model is relatively low. EAGLE/EAGLE3 follows and provides modest gains on most workloads. $n$-gram achieves its best speedups on highly repetitive code-editing tasks, where rare but very long bursts occur. Outside such settings, its heavy-tailed and unstable acceptance makes performance less predictable.




\begin{figure*}[t]
\centering
\lstset{style=codebox, language=Python}

\begin{minipage}[t]{0.485\textwidth}
\textbf{Prompt}
\begin{lstlisting}
class Car:

    def __init__(self, make, model, year, color):
        self.make = make
        self.model = model
        self.year = year
        self.color = color


car1 = Car('Toyota', 'Camry', 2018, 'Red')
print(car1.color)

# Replace the string data type used for car color with an enum data type. 
# Just output the code, do not include any explanation.
\end{lstlisting}
\end{minipage}\hfill
\begin{minipage}[t]{0.485\textwidth}
\textbf{Output}
\begin{lstlisting}
from enum import Enum


class Color(Enum):
    RED = 1
    BLUE = 2
    GREEN = 3
    YELLOW = 4
    BLACK = 5
    WHITE = 6


class Car:

    def __init__(self, make, model, year, color):
        self.make = make
        self.model = model
        self.year = year
        self.color = color


car1 = Car('Toyota', 'Camry', 2018, Color.RED)
print(car1.color)
\end{lstlisting}
\end{minipage}

\vspace{-0.4em}
\caption{\textbf{Example of prompt-output repetition in code editing from InstructCoder.} The output preserves most of the prompt and applies a localized edit, a structure that $n$-gram can exploit to produce occasional long accepted spans.}
\label{fig:ngram_prompt_output_full}
\vspace{-1.0em}
\end{figure*}
\begin{figure*}[h!]
\begin{center}
\captionsetup{justification=centering}
\begin{subfigure}{0.24\linewidth}
\includegraphics[width=\linewidth]{figures/BLEU_Score/instructcoder_meta-llama_Llama-3.1-8B-Instruct_ngram_vs_eagle_relative_heatmap.pdf}
\caption{Llama3.1-8B, $n$-gram-fixed-3 vs EAGLE}
\end{subfigure}
\begin{subfigure}{0.24\linewidth}
\includegraphics[width=\linewidth]{figures/BLEU_Score/instructcoder_meta-llama_Llama-3.1-8B-Instruct_ngram_vs_eagle3_relative_heatmap.pdf}
\caption{Llama3.1-8B, $n$-gram-fixed-3 vs EAGLE-3}
\end{subfigure}
\begin{subfigure}{0.24\linewidth}
\includegraphics[width=\linewidth]{figures/BLEU_Score/instructcoder_Qwen_Qwen3-8B_ngram_vs_eagle3_relative_heatmap.pdf}
\caption{Qwen3-8B, $n$-gram-fixed-3 vs EAGLE-3}
\end{subfigure}
\begin{subfigure}{0.24\linewidth}
\includegraphics[width=\linewidth]{figures/BLEU_Score/instructcoder_meta-llama_Meta-Llama-3-70B-Instruct_ngram_vs_eagle_relative_heatmap.pdf}
\caption{Llama3-70B, $n$-gram-fixed-3 vs EAGLE}
\end{subfigure}

% NRGAM-FIXED-5 Below
\begin{subfigure}{0.24\linewidth}
\includegraphics[width=\linewidth]{figures/BLEU_Score/ngram5-reduced-batchsize/instructcoder_meta-llama_Llama-3.1-8B-Instruct_ngram_vs_eagle_relative_heatmap.pdf}
\caption{Llama3.1-8B, $n$-gram-fixed-5 vs EAGLE}
\end{subfigure}
\begin{subfigure}{0.24\linewidth}
\includegraphics[width=\linewidth]{figures/BLEU_Score/ngram5-reduced-batchsize/instructcoder_meta-llama_Llama-3.1-8B-Instruct_ngram_vs_eagle3_relative_heatmap.pdf}
\caption{Llama3.1-8B, $n$-gram-fixed-5 vs EAGLE-3}
\end{subfigure}
\begin{subfigure}{0.24\linewidth}
\includegraphics[width=\linewidth]{figures/BLEU_Score/ngram5-reduced-batchsize/instructcoder_Qwen_Qwen3-8B_ngram_vs_eagle3_relative_heatmap.pdf}
\caption{Qwen3-8B, $n$-gram-fixed-5 vs EAGLE-3}
\end{subfigure}
\begin{subfigure}{0.24\linewidth}
\includegraphics[width=\linewidth]{figures/BLEU_Score/ngram5-reduced-batchsize/instructcoder_meta-llama_Meta-Llama-3-70B-Instruct_ngram_vs_eagle_relative_heatmap.pdf}
\caption{Llama3-70B, $n$-gram-fixed-5 vs EAGLE}
\end{subfigure}

\caption{Correlation between BLEU-4 and $n$-gram speedup on InstructCoder. Each heatmap shows $n$-gram’s relative speedup (\%) over EAGLE or EAGLE-3, grouped by BLEU-4 scores (x-axis) and batch sizes (y-axis). Blue means that $n$-gram outperforms the draft-model method, while red means the opposite. $n{=}$XYZ in the x-axis labels denotes the number of requests in each BLEU bucket. Figures (a) to (d) show the results when $n$-gram proposes 3 tokens per step and Figures (e) to (h) show the results when that is set to 5.}
\label{fig:ngram_bleu_heatmap}
\vspace{-0.2in}
\end{center}
\end{figure*}


%% <<< end acceptance_behavior.tex


%% >>> begin case_study.tex
\section{Case Study: $n$-gram Performance on InstructCoder}
\label{subsect:ngram-bleu-score}


The $n$-gram SD is especially attractive for speculative decoding because it is training-free.
We observe that $n$-gram achieves particularly strong speedups on the code-editing workload, such as InstructCoder, and we therefore conduct a focused case study to understand the underlying reasons.

We hypothesize that its advantage stems from the \emph{local repetition} inherent in code-editing tasks: when upcoming tokens have already appeared in the prompt, $n$-gram lookup is more likely to propose correct continuations that are subsequently accepted by the target model during verification.
To test this hypothesis, we quantify prompt–output overlap using BLEU-$n$~\cite{papineni2002bleu}.
BLEU-$n$ measures the fraction of overlapping $n$-grams, with higher values indicating stronger local reuse or copying.
For each request, we compute the BLEU score between the prompt and the output generated without speculative decoding, and then group them into 5 intervals: requests with BLEU score of 0-0.2, 0.2-0.4, 0.4-0.6, 0.6-0.8 and 0.8-1.0.
Within each interval, we compare the speedups achieved by $n$-gram and by EAGLE/EAGLE-3.
We report the results using BLEU-4 score by default~\cite{papineni2002bleu}. In our experiments, BLEU-4 gives the clearest speedup heatmap and shows a clear cutoff where, above a certain overlap range, $n$-gram beats EAGLE/EAGLE-3 for every batch size. We have also repeated the analysis with BLEU-1 through BLEU-10 scores, and observe the same overall patterns.

\begin{quoting}[leftmargin=1.5em, rightmargin=1.5em, vskip=0pt]
\textbf{Overlap Drives $n$-gram Speedups.}
\emph{$n$-gram works best on code editing because the output often reuses short spans that already appear in the prompt. As overlap increases, $n$-gram speedup rises and eventually surpasses EAGLE/EAGLE-3.}
\end{quoting}
We observe that higher BLEU-$n$ scores strongly correlate with greater $n$-gram speedup (\autoref{fig:ngram_bleu_heatmap}). 
For requests with lower BLEU-$n$ scores (little overlap), $n$-gram underperforms due to inaccurate proposals. Once the BLEU-$n$ score exceeds a certain threshold, (e.g. larger than $0.6$ for Llama3.1-8B), $n$-gram consistently outperforms EAGLE and EAGLE-3 across all batch sizes. With a proposal length of 3, it achieves up to $53\%$ higher speedup than EAGLE and EAGLE-3. When we further increase the proposal length of $n$-gram to 5, the boundary where $n$-gram performs better is clearer, and this brings an even higher performance of up to $100\%$ higher speedup than EAGLE/EAGLE-3.
This is likely because a larger proposal length allows $n$-gram to reuse longer repeated spans in the prompt or recent context when the overlap is high. 
Overall, these results confirm that $n$-gram speculation benefits directly from prompt-level repetition in code-editing workloads. 

Finally, our BLEU analysis measures overlap only between the complete prompt and the complete generated output. In practice, $n$-gram can match and reuse tokens from the full in-context history, which includes both the prompt and the tokens generated so far. We leave an analysis that measures overlap with respect to this evolving context for future work.

% A consistent trend emerges for non-reasoning models (i.e. Llama3.1-8B, Llama3-70B) and non-reasoning datasets (i.e. InstructCoder, ShareGPT and CNNDailyMail): as BLEU increases, the speedup of $n$-gram speculation rises and eventually surpasses that of the Eagle/Eagle-3 methods. Figure~X (Llama3.1-8B and Llama3-70B on \texttt{InstructCoder}) illustrates this relationship, confirming that greater prompt--output overlap directly results in better $n$-gram efficiency. \jerry{Add Bar chart here!} 

% However, this does not hold true for reasoning models (i.e. Qwen3-8B) and reasoning datasets (i.e. AIME22-24, GPQA-MAIN). We found that the similarity between the prompt and generated outputs are often too low for 

% \paragraph{Cutoffs where $n$-gram beats Eagle/Eagle-3.}
% To quantify when $n$-gram speculation becomes more effective than Eagle/Eagle-3 methods, we define a \emph{BLEU score cutoff} for each (model, dataset) pair. The cutoff is the smallest BLEU score bucket where the speedup of $n$-gram exceeds that of the compared Eagle or Eagle-3 method and remains higher in the subsequent bucket. We sweep and compute the BLEU-1 to BLEU-10 scores for every request across all configurations and identify the most appropriate cutoffs among them.
% \autoref{tab:ngram-cutoffs} summarizes the observed cutoffs across all configurations.

% \begin{table}[t]
% \centering
% \caption{BLEU score cutoffs where $n$-gram exceeds other speculative decoding (SD) methods.
% “Cmp.” denotes the SD method being compared. “Eg/Eg-3” refers to Eagle/Eagle-3.
% “Weak” indicates that fewer than 10 requests fall into BLEU buckets higher than the cutoff.
% "NF" means no reliable cutoff found because average speedups do not cross. 
% "NF-D" means majority of the requests have low BLEU score, "NF-C" means no reliable cutoff found despite that requests have a range of BLEU scores.
% \jerry{Double check the results for Qwen3-8B.}}
% \label{tab:ngram-cutoffs}
% \resizebox{\linewidth}{!}{
% \begin{tabular}{l l c c}
% \toprule
% \textbf{Model} & \textbf{Dataset} & \textbf{Cmp.} & \textbf{Cutoff} \\
% \midrule
% Llama3.1-8B & InstructCoder & Eg/Eg-3   & BLEU-4 $> 0.5$ \\
%             & ShareGPT      & Eg/Eg-3   & BLEU-4 $> 0.6$ (weak) \\
%             & CNNDailyMail  & Eg/Eg-3   & BLEU-4 $> 0.4$ (weak) \\
%             & GSM8K         & Eg/Eg-3 & NF-D(?) \\
% \midrule
% Llama3-70B  & InstructCoder & Eg   & BLEU-4 $> 0.6$ \\
%             & ShareGPT      & Eg   & BLEU-3 $> 0.7$ (weak) \\
%             & CNNDailyMail  & Eg   & NF-C \\
%             & GSM8K         & Eg   & NF-C \\
% \midrule
% Qwen3-8B    & InstructCoder & Eg-3   &  BLEU-3 $> 0.35$ (weak)\\
%             & ShareGPT      & Eg-3   &  NF-C\\
%             & CNNDailyMail  & Eg-3   &  NF-C\\
%             & GSM8K         & Eg-3   &  NF-D\\
%             & AIME21-24     & Eg-3   &  NF-D\\
%             & GPQA-Main     & Eg-3   &  NF-D\\
% \bottomrule
% \end{tabular}
% }
% \end{table}


%% <<< end case_study.tex


%% >>> begin optimal_sd.tex
\begin{figure*}[t]
  \centering
\begin{subfigure}[t]{0.35\linewidth}
\includegraphics[width=\linewidth]{figures/optimal_sd/oracle/legend_only.pdf}
\end{subfigure}

\begin{subfigure}[b]{0.22\linewidth}
    \includegraphics[width=\linewidth]{figures/optimal_sd/oracle/ngram_instructcoder_speedup.pdf}
    \caption{$n$-gram, InstructCoder.}
    \label{fig:speedup-oracle-ngram-instructcoder}
\end{subfigure}
\begin{subfigure}[b]{0.22\linewidth}
    \includegraphics[width=\linewidth]{figures/optimal_sd/oracle/ngram_sharegpt_speedup.pdf}
    \caption{$n$-gram, ShareGPT.}
    \label{fig:speedup-oracle-ngram-sharegpt}
\end{subfigure}
\begin{subfigure}[b]{0.22\linewidth}
    \includegraphics[width=\linewidth]{figures/optimal_sd/oracle/ngram_gsm8k_speedup.pdf}
    \caption{$n$-gram, GSM8K.}
    \label{fig:speedup-oracle-ngram-gsm8k}
\end{subfigure}
\begin{subfigure}[b]{0.22\linewidth}
    \includegraphics[width=\linewidth]{figures/optimal_sd/oracle/ngram_cnn_speedup.pdf}
    \caption{$n$-gram, CNN/DailyMail.}
    \label{fig:speedup-oracle-ngram-cnn}
\end{subfigure}

\begin{subfigure}[b]{0.22\linewidth}
    \includegraphics[width=\linewidth]{figures/optimal_sd/oracle/eagle_instructcoder_speedup.pdf}
    \caption{EAGLE, InstructCoder.}
    \label{fig:speedup-oracle-eagle-instructcoder}
\end{subfigure}
\begin{subfigure}[b]{0.22\linewidth}
    \includegraphics[width=\linewidth]{figures/optimal_sd/oracle/eagle_sharegpt_speedup.pdf}
    \caption{EAGLE, ShareGPT.}
    \label{fig:speedup-oracle-eagle-sharegpt}
\end{subfigure}
\begin{subfigure}[b]{0.22\linewidth}
    \includegraphics[width=\linewidth]{figures/optimal_sd/oracle/eagle_gsm8k_speedup.pdf}
    \caption{EAGLE, GSM8K.}
    \label{fig:speedup-oracle-eagle-gsm8k}
\end{subfigure}
\begin{subfigure}[b]{0.22\linewidth}
    \includegraphics[width=\linewidth]{figures/optimal_sd/oracle/eagle_cnn_speedup.pdf}
    \caption{EAGLE, CNN/DailyMail.}
    \label{fig:speedup-oracle-eagle-cnn}
\end{subfigure}
    
    % \vspace{-0.2in}
    \caption{Oracle vs fixed proposed length speedup on Llama3.1-8B.}
    \label{fig:speedup-oracle}
    % \vspace{-0.1in}
%
\end{figure*}


\begin{figure*}[h]
  \centering
      \begin{subfigure}[b]{0.24\linewidth}
    \includegraphics[width=\linewidth]{figures/optimal_sd/diff/diff_eagle3_ngram_instructcoder_acceptance_stats.pdf}
    \caption{InstructCoder}
    \label{fig:acc-diff-instrcutcoder}
  \end{subfigure}
      \begin{subfigure}[b]{0.24\linewidth}
    \includegraphics[width=\linewidth]{figures/optimal_sd/diff/diff_eagle3_ngram_sharegpt_acceptance_stats.pdf}
    \caption{ShareGPT}
    \label{fig:acc-diff-sharegpt}
  \end{subfigure}
      \begin{subfigure}[b]{0.24\linewidth}
    \includegraphics[width=\linewidth]{figures/optimal_sd/diff/diff_eagle3_ngram_gsm8k_acceptance_stats.pdf}
    \caption{GSM8K}
    \label{fig:acc-diff-gsm8k}
  \end{subfigure}
          \begin{subfigure}[b]{0.24\linewidth}
    \includegraphics[width=\linewidth]{figures/optimal_sd/diff/diff_eagle3_ngram_cnn_acceptance_stats.pdf}
    \caption{CNN/DailyMail}
    \label{fig:acc-diff-cnn}
  \end{subfigure}
  \caption{Per-position accepted-length difference between $n$-gram and EAGLE on Llama 3.1-8B. Red indicates positions where EAGLE accepts longer spans, while blue indicates positions favoring $n$-gram. The figure illustrates that different speculative decoding methods exhibit distinct acceptance behaviors across decoding positions.}
\label{fig:acc-diff}
\end{figure*}


\begin{figure*}
\centering
\begin{subfigure}[t]{0.65\linewidth}
\includegraphics[width=\linewidth]{figures/optimal_sd/oracle/combined_legend_only.pdf}
\end{subfigure}

    \begin{subfigure}[b]{0.22\linewidth}
      \includegraphics[width=\linewidth]{figures/optimal_sd/oracle/combined_instructcoder_speedup.pdf}
      \caption{InstructCoder}
      \label{fig:opt-speedup-instrcutcoder}
  \end{subfigure}
  \begin{subfigure}[b]{0.22\linewidth}
      \includegraphics[width=\linewidth]{figures/optimal_sd/oracle/combined_sharegpt_speedup.pdf}
      \caption{ShareGPT}
      \label{fig:opt-speedup-sharegpt}
  \end{subfigure}
   \begin{subfigure}[b]{0.22\linewidth}
      \includegraphics[width=\linewidth]{figures/optimal_sd/oracle/combined_gsm8k_speedup.pdf}
      \caption{GSM8K}
      \label{fig:opt-speedup-gsm8k}
  \end{subfigure}
      \begin{subfigure}[b]{0.22\linewidth}
      \includegraphics[width=\linewidth]{figures/optimal_sd/oracle/combined_cnn_speedup.pdf}
      \caption{CNN/DailyMail}
      \label{fig:opt-speedup-cnn}
  \end{subfigure}
  \caption{Combining different SD methods for optimal speedup on Llama3.1-8B.}
  \label{fig:opt-speedup}
\end{figure*}



\section{Theoretical Speedup Upper Bound of Speculative Decoding}
\label{sec:theoretical-speedup}
% With the understanding of acceptance behavior, we can now estimate the potential speedup if we could perfectly predict the number of accepted tokens at each position. To estimate zero-waste speculative decoding, we proposed generating a very long speculative sequence—for example, 20 tokens at each position—to observe how many tokens could be accepted at each step. After each generation step, even if multiple tokens were accepted, we discarded all but one and advanced to the next position. Essentially, we recorded the maximum number of tokens that could be accepted at each speculative decoding step. Using this information, along with the execution time breakdown from the previous section, we can estimate the potential speedup if we knew the exact number of accepted tokens at each position in advance.
In this section, we aim to derive an upper bound of the speedup achievable by SD. Our goal is to understand how far current SD methods are from the theoretical optimum, assess their current performance limits, and highlight potential directions for future research.

As discussed in \autoref{sec:breakdown}, verification time dominates the overall execution cost. This naturally raises the question: \emph{if we could eliminate wasted verification by avoiding the verification of tokens that would ultimately be rejected, what is the maximum speedup that speculative decoding could achieve?
In other words, assuming an oracle—or a sufficiently accurate mechanism—that proposes only tokens likely to be accepted, what is the upper bound on the achievable speedup?} In the following section, we try to derive this number and understand the gap between current speedup and theoritical upper bound.

Moreover, as observed in \autoref{sec:acc-behavior}, different SD variants exhibit distinct acceptance behaviors. Specifically, they differ in the number of tokens accepted at each generation step. To better understand the true performance upper bound, we investigate whether combining different SD strategies could yield an even more optimal acceptance pattern and, consequently, higher overall efficiency. 

%therefore conclude with a list of open challenges for future work inspired by our study.



\subsection{Minimal Verification Cost}
We begin by addressing the following question: if we could minimize the verification cost, what is the maximum speedup that speculative decoding could achieve? To explore this, we compare the performance of using a fixed proposed length against that of an oracle-based proposed length. The oracle setup assumes that the number of tokens accepted at each generation step is known in advance. At each step, we set the proposed length equal to the actual accepted length, ensuring that all proposed tokens are accepted. This setup effectively simulates the ideal case where verification incurs no rejection overhead.

One might argue that if all proposed tokens can be perfectly accepted, the large model would no longer be necessary. While that is theoretically correct, it is practically unattainable—no proposal mechanism can perfectly predict future target tokens. Hence, the oracle configuration serves purely as an upper bound on the achievable speedup of speculative decoding, highlighting the performance ceiling that real methods can approximate.

We conduct this analysis using Llama 3.1-8B on the InstructorCoder, ShareGPT, GSM8K, and CNN/DailyMail datasets. We observe a substantial gap between the oracle speedup and the speedup achieved with a fixed proposal length. The oracle speedup represents an upper bound in which the accepted length at each decoding step is known in advance, and the system proposes exactly that many tokens. As a result, every proposed token is accepted, eliminating any wasted speculation.

 As shown in \autoref{fig:speedup-oracle}, on the Instructcoder dataset with a batch size of one and $n$-gram as the SD method, the oracle setup achieves a speedup of approximately 2.75$\times$, whereas the fixed proposed length configuration attains about 2.1$\times$ for the best proposed length of 5. Moreover, the gap generally widens as batch size increases: the fixed-$k$ curves drop much faster, while the oracle speedup degrades more gently. This trend is consistent with our observation in \autoref{fig:exec-breakdown}. As batch size grows, verification takes up a larger share of the decoding time, so fixed-$k$ methods would incur larger overhead for verifying draft tokens that are ultimately rejected.

% An interesting pattern emerges when comparing $n$-gram and EAGLE. With small batch sizes, the gap between oracle and fixed-length speedups is noticeably smaller on ShareGPT. As discussed in \autoref{sec:acc-behavior}, when $n$-gram successfully identifies a match, it tends to be correct, leading to a higher acceptance rate—that is, fewer rejected tokens and reduced verification overhead. Consequently, there is less room for further reducing verification cost. However, as the batch size increases, the system becomes increasingly compute-bound, causing the performance gap to widen for both methods.



% \begin{figure}[t]
%   \centering
%   \begin{subfigure}[b]{0.45\linewidth}
%     \includegraphics[width=\linewidth]{figures/optimal_sd/oracle/eagle3_instructcoder_speedup.pdf}
%     \caption{Instructcoder, Eagle}
%   \end{subfigure}
%   \begin{subfigure}[b]{0.45\linewidth}
%     \includegraphics[width=\linewidth]{figures/optimal_sd/oracle/eagle3_sharegpt_speedup.pdf}
%     \caption{ShareGPT, Eagle}
%   \end{subfigure}
%   \caption{Oracle, preliminary predictor, and fixed proposed length speedup for Eagle.}
%   \label{fig:speedup-oracle-eagle}
% \end{figure}

% \textbf{Heuristic-based Accepted Length Predictor.} While previous work~\cite{agrawal2024adaedl, zhang2025draftmodelknowsstop} explores different means to reduce verification cost, we instead aim to understand how the current optimization methods compare to the theoretical upper bound and whether simple heuristics can help close this gap.
% We experiment with a lightweight predictor that leverages the draft method’s confidence to decide whether to send proposed tokens for verification. In SD, after the draft method generates a sequence of candidate tokens, only those with confidence above a threshold (e.g., 0.3 or 0.7) are forwarded for verification, while low-confidence tokens are skipped. The detailed algorithm is described in Appendix~\ref{appendix:predictor}.
% As shown by the purple lines in \jerry{todo}, the confidence-based predictor yields modest improvements in the speedup and narrows the gap toward the oracle performance. However, the remaining gap is still substantial. The predictor must be both lightweight—since it lies on the critical path—and sufficiently accurate to effectively reduce unnecessary verification of rejected tokens. We believe this direction holds great potential for further improving the efficiency of speculative decoding.


\subsection{An Adaptive Approach to Achieve Optimal Speedup}
Although EAGLE generally outperforms the $n$-gram method, 
% —for example, on the ShareGPT dataset, the oracle speedup of $n$-gram peaks at only about 1.2$\times$, which remains below EAGLE’s performance—
a closer inspection of acceptance behavior reveals complementary strengths between the two. 
As illustrated in \autoref{fig:acc-diff}, the per-position differences in accepted tokens exhibit alternating regions of red (favoring EAGLE) and blue (favoring $n$-gram), indicating that the two methods excel under different conditions. This behavior is intuitive: for example, in code snippets with strong local regularities or recurring grammatical patterns, $n$-gram can often find exact matches, which are likely to be correct.
In such cases, because $n$-gram incurs substantially lower proposing overhead, it can be the more efficient choice. At other positions where local repetition is weaker, EAGLE serves as a robust fallback, since it proposes correct tokens across diverse contexts more consistently.



This naturally brings up the next question: \emph{can we combine their advantages to achieve even higher overall speedups?}

To explore the upper bound on achievable speedup, we assume a \emph{perfect} predictor that can (1) determine which method performs best at each position, and (2) accurately predict the number of tokens that will be accepted at that position. Under this idealized assumption, we measure the maximum possible speedup attainable by speculative decoding.

It is important to note the simplifications in this setup, which represent an upper bound on achievable performance. First, we assume the predictor is flawless in determining the better method for each position. Second, we presume perfect knowledge of the number of tokens each method can accept at every position. 
% Third, we simplify the overhead comparison by using $n$-gram's proposing overhead, which is generally lower than Eagle's. This choice allows us to estimate the best-case scenario. 
Finally, we do not account for the overhead of maintaining the KV cache for EAGLE's proposing head, which would otherwise require partial prefills if a request was paused and resumed.


Under these assumptions, we report the best achievable speedup in \autoref{fig:opt-speedup}. \textbf{Oracle\_Combine} represents a theoretical upper bound in which, at each generation position, the method (EAGLE or $n$-gram) that yields the longer accepted span is selected. The resulting accepted length is then used as the proposed length. Compared to the best fixed strategy, \textbf{Oracle\_Combine} exposes substantial additional headroom, achieving up to a $2.2\times$ further speedup.

The size of this headroom depends strongly on the workload. InstructCoder shows the largest gap between existing methods and Oracle\_Combine, matching \autoref{fig:acc-diff} where red and blue regions alternate frequently, i.e., there are many positions where one method clearly outperforms the other. ShareGPT and CNN/DailyMail exhibit similar but smaller gains, indicating more limited but still meaningful room for the methods to complement each other. In contrast, GSM8K offers little additional benefit from combining because $n$-gram rarely achieves long acceptances, so \textbf{Oracle\_Combine} stays close to \textbf{Oracle\_Eagle}. 

Taken together, these results point to a promising direction for future work: developing an accurate yet lightweight predictor capable of adapting to varying levels of acceptance behavior across workloads, requests and token positions.

% Let the figures stay with this section, instead of overflowing to the reference page.
% \clearpage

% With these assumptions, we present the optimal speedup by combining the two methods in \autoref{fig:opt-speedup}. Specifically, \textbf{Oracle\_Combine} represents the theoretical upper bound obtained by selecting, at each position, the SD variant (EAGLE or $n$-gram) that yields the longer accepted sequence. 
% We observe a substantial gap between the current methods and oracle speedups, reaching over 3$\times$ improvement when the batch size is below 32. This indicates a promising direction for future work, though it poses challenges such as developing an accurate yet lightweight predictor capable of adapting to varying levels of acceptance behavior variance.

% \begin{figure}[h]
%   \centering
%   \begin{subfigure}[b]{0.35\linewidth}
%     \includegraphics[width=\linewidth]{figures/optimal_sd/oracle/combined_instructcoder_speedup.pdf}
%     \caption{InstructCoder}
%   \end{subfigure}
%   \begin{subfigure}[b]{0.6\linewidth}
%     \includegraphics[width=\linewidth]{figures/optimal_sd/oracle/singlepage_plot_plus_legend.pdf}
%     \caption{ShareGPT}
%    \end{subfigure}
%   \caption{Oracle combined speedup.}
%   \label{fig:acc-combine}
% \end{figure}
%% <<< end optimal_sd.tex


%% >>> begin conclusion.tex
\section{Conclusion}
This work presents the first systematic evaluation of speculative decoding (SD) in a real, optimized inference engine. By systematically dissecting end-to-end performance across workloads and SD variants, we identify verification as the dominant bottleneck and reveal strong variability in acceptance behavior across positions, requests, and datasets. Leveraging these insights, we quantify the theoretical upper bound of SD speedup and expose the gap between observed and ideal efficiency. Together, these findings deepen our understanding of SD’s practical behavior and highlight opportunities that can further unlock its full potential in large-scale inference systems.

\section{Acknowledgements}
We are grateful to Tomas Ruiz for implementing draft-model-based speculative decoding in vLLM, which allowed us to profile and compare this approach. We also thank Yilong Zhao, Yifan Qiao, and other members of the Sky Computing Lab for their helpful discussions and feedback. This work was supported in part by a gift from NVIDIA, including the DGX server used in this study. Jiaxiang is supported by the Singapore National Science Scholarship and the NUS Development Grant. Any opinions, findings, and conclusions or recommendations expressed in this paper are those of the authors and do not necessarily reflect the views of the supporting organizations.
%% <<< end conclusion.tex

\bibliography{example_paper}
\bibliographystyle{mlsys2025}


%% >>> begin appendix.tex
\clearpage
\appendix
\section{Appendix}

\subsection{Dataset}
\label{appendix:dataset}
% \textbf{CNNDailyMail}
% \textbf{ShareGPT}
% \textbf{InstructCoder}
% \textbf{GSM8K}
% \textbf{AIME22-24}
% \textbf{GPQA-Main}
\begin{figure}[h]
    \centering
    \includegraphics[width=\linewidth]{figures/workload_length_distribution/before_filtering/all_models_combined.pdf}
    \caption{Dataset Input/Output Length distributions.}
    \label{fig:dataset-input-output-plot}
\end{figure}

\subsection{Memory Breakdown Calculations}
\label{appendix:mem}
\paragraph{Static memory (GiB).} Assuming FP16 weights (2 bytes/param), the static GPU memory for model weights:
\[
M_\text{static,GiB}=\frac{(P_\text{target}+P_\text{draft})\cdot 2}{2^{30}}
\]
Here, $P_\text{target}$ is the target model's parameter size and $P_\text{draft}$ includes any additional parameters introduced by the speculative decoding method, both in units of Billions.

For EAGLE/EAGLE3, $P_\text{draft}$ refers to the parameters used in speculative decoding. For the EAGLE models we used, $P_\text{draft}$ refers to the parameter size of the autoregressive head, which includes one decoding layer and one fully connected (FC) layer. They share the same embedding layer and language modeling (LM) head with the target model. For the EAGLE3 models we used, $P_\text{draft}$ refers to the parameter size of the entire EAGLE3 model. It includes one decoding layer, the multi-layer feature fusion FC layer, a final normalization layer and its own LM head. In addition, the model checkpoints include two 1-dimensional token-id remapping tables (\texttt{t2d} and \texttt{d2t}) that translate between the target model’s vocabulary and the draft LM head’s vocabulary; they typically contribute negligibly to the parameter memory. When loading model weights in vLLM v0.10.1.1, the \texttt{t2d} tensors are skipped, and we exclude them when counting the parameter size. Similarly, these EAGLE3 models reuse the same embedding layer from the target models.

For Qwen3-0.6B, \textit{tie\_word\_embeddings} defaults to True and its LM head shares the same weight matrix as the embedding layer; thus, we exclude the size of its language modeling (LM) head.

For example, the static memory for Llama3-70B-Instruct with Llama3.2-1B-Instruct as the draft model is $(70.55+1.23)*10^9*2/2^{30}=~133.7GiB$. For Llama3.1-8B-Instruct with EAGLE-LLaMA3.1-Instruct-8B, it is $(8.03+0.25)*10^9*2/2^{30}=15.42GiB$.

\paragraph{Per-token KV cache (KiB).}
To calculate the KV cache size of each generated token:
\[
M_\text{KV/token,KiB}=
\frac{L_\text{h}\cdot 2 \cdot n_\text{kv}\cdot d_\text{head}\cdot 2}{2^{10}} .
\]
$L_\text{h}$ is the number of hidden layers,
$n_\text{kv}$ is the number of KV heads (since our models are all Grouped-Query-Attention-based),
$d_\text{head}$ is the head dimension, the first factor $2$ is for both Key and Value, and the last factor $2$ is bytes/element for FP16. For Llama3-70B-Instruct without speculative decoding or model-free SD methods like $n$-gram, it is $80*2*8*128*2/2^{10}=~320KiB$. Similarly, for Llama3.1-8B-Instruct, it is $32*2*8*128*2/2^{10}=~128KiB$.

For draft-model-based SD, the total per-token KV cache is the sum over each generated token by the target and draft model. For Llama3-70B-Instruct with Llama3.2-1B-Instruct as the draft model, the per-token KV cache memory size is $80*2*8*128*2/2^{10}$ + $16*2*8*64*2/2^{10}$ = $352KiB$.

Similarly, for EAGLE-based SD, the total per-token KV cache is the sum over each generated token by the target and EAGLE heads. For Llama3-70B-Instruct with Llama3.2-1B-Instruct as the draft model, the per-token KV cache memory size is $80*2*8*128*2/2^{10}$ + $1*2*8*128*2/2^{10}$ = $324KiB$. The value 1 here means that EAGLE uses an additional transformer block for proposal generation, which effectively adds one extra layer.

\begin{table}[t]
\centering
\small
\setlength{\tabcolsep}{4pt}
\begin{tabular}{l|rrrr}
\toprule
Model & $P$ & $L$ & $n_\text{kv}$ & $d_\text{head}$ \\
\midrule
Llama3.1-8B-Instruct  & 8.03B & 32 & 8 & 128 \\
Llama3-70B-Instruct   & 70.55B & 80 & 8 & 128 \\
Llama3.2-1B-Instruct  & 1.23B & 16 & 8 & 64 \\
Qwen3-8B              & 8.19B & 36 & 8 & 128 \\
Qwen3-0.6B              & 0.596B & 28 & 8 & 128 \\
EAGLE-LLaMA3.1-Instruct-8B & 0.25B & 1 & 8 & 128 \\
EAGLE3-LLaMA3.1-Instruct-8B & 0.425B & 1 & 8 & 128 \\
EAGLE-LLaMA3-Instruct-70B & 0.99B & 1 & 8 & 128 \\
EAGLE3-Qwen3-8B            &  0.40B     & 1 & 8 & 128 \\
\bottomrule
\end{tabular}
\vspace{2pt}
\caption{Model specs used for memory calculations. $P$ is the parameter size of the weights used, obtained by counting their parameters; $L$ is the number of hidden layers, $n_\text{kv}$ is the number of Key/Value heads, and $d_\text{head}$ is the head dimension, obtained from the model's configuration files on Hugging Face.}
\label{tab:mem_model_specs}
\vspace{-0.10in}
\end{table}


\subsection{Acceptance Behaviour}
\label{appendix:acc-behaviour}

% \begin{figure}[h!]
%   \centering
%  \begin{subfigure}{0.45\linewidth}
%     \includegraphics[width=\linewidth]{figures/acc_len/heatmap/instructcoder_eagle_llama3.1-8B_acceptance_stats.pdf}
%     \caption{InstructCoder, EAGLE}
%   \end{subfigure}
%   \hspace{1pt}
%   \begin{subfigure}{0.45\linewidth}
%     \includegraphics[width=\linewidth]{figures/acc_len/heatmap/sharegpt_eagle_llama3.1-8B_acceptance_stats.pdf}
%     \caption{ShareGPT, EAGLE}
%   \end{subfigure}
%   \caption{Generated length per token position for Llama3.1-8B. Requests are sorted based on generation length.}
%   \label{fig:acceptance-length-heatmap-continued}
% \end{figure}

% \begin{figure} [h!]
% \captionsetup{justification=centering}
%   \centering
%     \hspace{1pt}
%     \begin{subfigure}{0.45\linewidth}
%     \includegraphics[width=\linewidth]{figures/acc_len/token_level/overlaid_binned_positions_gpqa_main_Qwen3-8B_bin1_range4190-7811.pdf}
%     \caption{Requests with outputs of 4K-8K tokens}
%   \end{subfigure}
%     \hspace{1pt}
%     \begin{subfigure}{0.45\linewidth}
%     \includegraphics[width=\linewidth]{figures/acc_len/token_level/overlaid_binned_positions_gpqa_main_Qwen3-8B_bin2_range7811-13351.pdf}
%     \caption{Requests with outputs of 8K-13K tokens}
%   \end{subfigure}
% \caption{Token-level: Average generated token length per generated token position for Qwen3-8B-Thinking on GPQA-Main.}
%   \label{fig:acc-len-token-level-continued}
% \end{figure}

\begin{figure}[h]
  \centering
  \includegraphics[width=0.7\linewidth]{figures/acc_len/request_level/legend.pdf}
  
    \begin{subfigure}{0.3\linewidth}
    \includegraphics[width=\linewidth]{figures/acc_len/request_level/box_by_model_aime_acc_len.pdf}
    \caption{AIME22-24}
  \end{subfigure}
      \hspace{5pt}
    \begin{subfigure}{0.3\linewidth}
    \includegraphics[width=\linewidth]{figures/acc_len/request_level//box_by_model_gpqa_main_acc_len.pdf}
    \caption{GPQA-Main}
  \end{subfigure}
\caption{Request-level generated length across datasets and models. Each box shows the distribution of per-request mean accepted length, with the box spanning the 25th–75th percentiles and whiskers covering the 5th–95th percentiles. Model used is Qwen3-8B-Thinking.}
  \label{fig:acc-len-request-level-box-plot-reasoning}
\end{figure}

\begin{table}[h!]
\centering
\begin{tabular}{llccc}
\toprule
Dataset & Method & Mean & Median & Std \\
\midrule
InstructCoder & $n$-gram & 7.27 & 6.57 & 4.19 \\
InstructCoder & EAGLE & 4.24 & 4.21 & 1.07 \\
InstructCoder & EAGLE-3 & 4.18 & 4.06 & 0.99 \\
CNN/DailyMail & EAGLE-3 & 3.08 & 2.97 & 0.61 \\
CNN/DailyMail & EAGLE & 2.32 & 2.31 & 0.21 \\
CNN/DailyMail & $n$-gram & 2.33 & 1.96 & 1.08 \\
ShareGPT & EAGLE-3 & 3.06 & 2.92 & 0.87 \\
ShareGPT & EAGLE & 2.82 & 2.68 & 0.80 \\
ShareGPT & $n$-gram & 2.02 & 1.37 & 1.87 \\
GSM8K & EAGLE-3 & 3.02 & 2.92 & 0.47 \\
GSM8K & EAGLE & 2.56 & 2.51 & 0.25 \\
GSM8K & $n$-gram & 1.41 & 1.28 & 0.45 \\
\midrule
GPQA\_Main & $n$-gram & 4.20 & 3.73 & 2.24 \\
GPQA\_Main & EAGLE-3 & 2.78 & 2.49 & 0.87 \\
AIME22-24 & $n$-gram & 3.53 & 2.96 & 2.29 \\
AIME22-24 & EAGLE-3 & 2.62 & 2.50 & 0.54 \\
\bottomrule
\end{tabular}
\caption{Dataset-level generated length by method, aggregated across models. For non-reasoning workloads, the average is aggregated over the results for Llama3.1-8B and Llama3-70B. For reasoning workloads, the average is aggregated over that of Qwen3-8B-Thinking.}
\label{tab:dataset_agg_by_method}
\end{table}

\clearpage
% \subsection{Heuristic-based Accepted Length Predictor}

% \label{appendix:predictor}
% \begin{algorithm}[h]
% \caption{Speculative Decoding with Confidence-Based Token Filtering}
% \begin{algorithmic}[1]
% \REQUIRE Draft model $D$, Target model $T$, confidence threshold $\theta$
% \STATE $x \gets$ initial context
% \WHILE{not end-of-sequence}
%     \STATE $Y \gets D(x)$ \COMMENT{Draft model proposes a sequence of tokens}
%     \STATE $Y' \gets [\,]$ \COMMENT{Tokens to be verified by the target model}
%     \FOR{each token $y_i$ in $Y$}
%         \STATE $c_i \gets \text{confidence}(y_i)$
%         \IF{$c_i \geq \theta$}
%             \STATE Append $y_i$ to $Y'$
%         \ELSE
%             \STATE \textbf{break} \COMMENT{Stop at first low-confidence token}
%         \ENDIF
%     \ENDFOR
%     \STATE $Z \gets T(x, Y')$ \COMMENT{Target model verifies selected tokens}
%     \STATE $x \gets x \circ \text{Accept}(Z)$ \COMMENT{Append accepted tokens to context}
% \ENDWHILE
% \end{algorithmic}
% \end{algorithm}

% \clearpage
% \subsection{End-to-End Experiment Setup}
\begin{table*}[t!]
\centering
\caption{Main experiment configurations across datasets, models, and speculative decoding (SD) methods used in End-to-End Speedup Measurement. DM stands for Draft-Model-Based SD.}
\label{tab:exp-config}
\setlength{\tabcolsep}{2.5pt}
\renewcommand{\arraystretch}{1.1}
\begin{tabular}{l l l l c c}
\toprule
\textbf{Dataset} & \textbf{Model} & \textbf{Hardware} & \textbf{SD Method} & \textbf{Max Output Len.} & \textbf{Batch Size}\\
\midrule
\texttt{CNNDailyMail} & Llama3.1-8B-Instruct & 1$\times$H100 & $n$-gram, EAGLE, EAGLE-3  & 8K & 1–128 \\
\texttt{ShareGPT}     & Llama3.1-8B-Instruct & 1$\times$H100 & $n$-gram, EAGLE, EAGLE-3  & 8K & 1–128 \\
\texttt{InstructCoder}& Llama3.1-8B-Instruct & 1$\times$H100 & $n$-gram, EAGLE, EAGLE-3  & 8K & 1–128 \\
\texttt{GSM8K}        & Llama3.1-8B-Instruct & 1$\times$H100 & $n$-gram, EAGLE, EAGLE-3  & 8K & 1–128 \\
\midrule
\texttt{CNNDailyMail} & Llama3-70B-Instruct  & 4$\times$H100 & $n$-gram, EAGLE, DM       & 8K & 1–128 \\
\texttt{ShareGPT}     & Llama3-70B-Instruct  & 4$\times$H100 & $n$-gram, EAGLE, DM         & 8K & 1–128 \\
\texttt{InstructCoder}& Llama3-70B-Instruct  & 4$\times$H100 & $n$-gram, EAGLE, DM         & 8K & 1–128 \\
\texttt{GSM8K}        & Llama3-70B-Instruct  & 4$\times$H100 & $n$-gram, EAGLE, DM         & 8K & 1–128 \\
\midrule
\texttt{CNNDailyMail} & Qwen3-8B    & 1$\times$H100 & $n$-gram, EAGLE-3, DM (v0.11.1rc1)         & 8K  & 1–128 \\
\texttt{ShareGPT}     & Qwen3-8B    & 1$\times$H100 & $n$-gram, EAGLE-3, DM (v0.11.1rc1)       & 8K  & 1–128 \\
\texttt{InstructCoder}& Qwen3-8B    & 1$\times$H100 & $n$-gram, EAGLE-3, DM (v0.11.1rc1)     & 8K  & 1–128 \\
\texttt{GSM8K}        & Qwen3-8B    & 1$\times$H100 & $n$-gram, EAGLE-3, DM (v0.11.1rc1)     & 8K  & 1–128 \\
\midrule
\texttt{CNNDailyMail} & Qwen3-8B-Thinking    & 1$\times$H100 & $n$-gram, EAGLE-3         & 8K  & 1–128 \\
\texttt{ShareGPT}     & Qwen3-8B-Thinking    & 1$\times$H100 & $n$-gram, EAGLE-3      & 8K  & 1–128 \\
\texttt{InstructCoder}& Qwen3-8B-Thinking    & 1$\times$H100 & $n$-gram, EAGLE-3    & 8K  & 1–128 \\
\texttt{GSM8K}        & Qwen3-8B-Thinking    & 1$\times$H100 & $n$-gram, EAGLE-3    & 8K  & 1–128 \\
\midrule
\multicolumn{6}{l}{\textit{Reasoning workloads}} \\
\texttt{AIME(22--24)} & Qwen3-8B-Thinking    & 1$\times$H100 & $n$-gram, EAGLE-3       & 32K & 1–16 \\
\texttt{GPQA-Main}    & Qwen3-8B-Thinking    & 1$\times$H100 & $n$-gram, EAGLE-3       & 32K & 1–16 \\
\texttt{AIME(22--24)} & GLM-4.5-Air    & 4$\times$H100 & $n$-gram, MTP (v0.11.1rc1)       & 32K & 1–4 \\
\texttt{GPQA-Main}    & GLM-4.5-Air     & 4$\times$H100 & $n$-gram, MTP (v0.11.1rc1)      & 32K & 1–4 \\
\bottomrule
\end{tabular}
\end{table*}

% \clearpage
% \subsection{Model Configurations}
\begin{table*}[t]
\centering
\caption{
Model configurations for EAGLE and Draft-Model-based SD.
Each configuration specifies the target model, the associated draft (or auxiliary) model from Hugging Face,
and the number of speculative tokens proposed per decoding step.}
\label{appendix:model-config}
\begin{tabular}{l l c l}
\toprule
\textbf{Variant} & \textbf{Target Model} & \textbf{\# Draft Tokens} & \textbf{Draft / Auxiliary Model (Hugging Face)} \\
\midrule
\textbf{EAGLE} 
& Llama3.1-8B-Instruct & 3 & \texttt{yuhuili/EAGLE-LLaMA3.1-Instruct-8B} \\
& Llama3-70B-Instruct  & 3 & \texttt{yuhuili/EAGLE-LLaMA3-Instruct-70B} \\
\midrule
\textbf{EAGLE-3} 
& Llama3.1-8B-Instruct & 3 & \texttt{yuhuili/EAGLE3-LLaMA3.1-Instruct-8B} \\
& Qwen3-8B             & 3 & \texttt{AngelSlim/Qwen3-8B\_eagle3} \\
\midrule
\textbf{Draft-Model-Based} 
& Llama3-70B-Instruct  & 3 & \texttt{meta-llama/Llama-3.2-1B-Instruct} \\
& Qwen3-8B             & 3 & \texttt{Qwen/Qwen3-0.6B} \\
\bottomrule
\end{tabular}
\end{table*}
\pagebreak


\pagebreak
\section{Artifact Appendix}

%%%%%%%%%%%%%%%%%%%%%%%%%%%%%%%%%%%%%%%%%%%%%%%%%%%%%%%%%%%%%%%%%%%%%
\subsection{Abstract}

This artifact contains two components.
The first is a vLLM-based profiling suite that includes implementation and benchmarking scripts for speculative decoding in vLLM, covering end-to-end throughput, time breakdown, and acceptance rate experiments across multiple models and methods (N-gram, EAGLE, EAGLE-3, Draft Model, MTP). Full reproduction of this suite requires Linux, Python 3.10, CUDA 12.8, conda/uv, HuggingFace access to gated Llama-3 models, a one-time ShareGPT download, and NVIDIA H100-80GB GPUs (1 GPU for 8B models and 4 GPUs for 70B/106B-scale models).
The second component is a lightweight simulator that runs on any laptop with Python 3.10 and uses included pre-profiled acceptance traces to reproduce the paper’s upper-bound and combined-proposer analyses.
Evaluators can validate the artifact by running the profiling scripts to regenerate the main end-to-end results for Llama3.1-8B, or/and by running the simulator to reproduce the oracle upper-bound analyses. Expected outputs are PDF figures whose trends closely match Figure 1(a)-(d) , Figure 9, and Figure 11.

\subsection{Artifact check-list (meta-information)}

{\small
\begin{itemize}
  \item {\bf Algorithm: } N-gram, EAGLE, EAGLE-3, Draft Model, MTP speculative decoding; simulation of speculative decoding performance.
  \item {\bf Program: } (1) Python~3.10, vLLM (built from source per branch). (2) Python~3.10, \texttt{tqdm}, \texttt{pandas}, \texttt{matplotlib}.
  \item {\bf Compilation: } (1) \texttt{uv pip install -e .} via \texttt{rebuild\_env.sh}; CUDA~12.8. (2) \texttt{pip install -r requirements.txt}; no GPU required.
  \item {\bf Data set: } (1) InstructCoder, CNN/DailyMail, GSM8K, AIME, GPQA (HuggingFace, auto-downloaded); ShareGPT (one-time download, see \S\ref{sec:ae-install}). (2) Pre-profiled acceptance data included in \texttt{simulator/data/}.
  \item {\bf Run-time environment: } (1) Linux, CUDA~12.8, conda. (2) Any OS with Python~3.10.
  \item {\bf Hardware: } (1) 1x-4x H100-80GB GPU (2) No GPU required.
  \item {\bf Run-time state: } (1) No other process should occupy the target GPU(s). (2) N/A.
  \item {\bf Execution: } (1) \texttt{run-l3-8b.sh} for Llama3.1-8B end-to-end performance profiling. (2) \texttt{./simulate\_and\_plot.sh}.
  \item {\bf Metrics: } (1) Throughput speedup, acceptance rate, time breakdown x; (2) simulated speedup vs.\ batch size.
  \item {\bf Output: } (1) \texttt{scripts/results/run\_<timestamp>/} with \texttt{.jsonl} results and logs. (2) \texttt{results/\{proposer\}\_\{dataset\}\_speedup.csv} and PDF figures in \texttt{figures/llama3.1-8B/}.
  \item {\bf Experiments: } (1). \texttt{run-l3-8b.sh} will take ~24-36 hours to complete on 1x H100. (2). \texttt{simulate\_and\_plot.sh} will take ~30-40 minutes to finish.
  \item {\bf How much disk space required (approximately)?: } recommended to be in 500GB of free disk space.
  \item {\bf How much time is needed to prepare workflow (approximately)?: } (1) $\approx$30-40\,min for each environment building. (2) $<$5\,min.
  \item {\bf How much time is needed to complete experiments (approximately)?: } See Experiments above.
  \item {\bf Publicly available?: } Yes.
  \item {\bf Data licenses (if publicly available)?: } See HuggingFace model/dataset cards if applicable.
\end{itemize}
}

%%%%%%%%%%%%%%%%%%%%%%%%%%%%%%%%%%%%%%%%%%%%%%%%%%%%%%%%%%%%%%%%%%%%%
\subsection{Description}
\label{sec:ae-desc}

\subsubsection{How to access}

(1) vLLM suite: \href{https://github.com/SpecDecode-Bench/vllm}{https://github.com/SpecDecode-Bench/vllm}; see \texttt{README.md} for more information.
(2) Simulator: \href{https://github.com/SpecDecode-Bench/simulator}{https://github.com/SpecDecode-Bench/simulator}; see \texttt{README.md} for more information.

\subsubsection{Hardware dependencies}

(1) Profiling: 1x H100-80GB GPU for 8B models; 4x H100-80GB GPUs for larger models.
(2) Simulation: No GPU required; any machine with Python~3.10 suffices. 
% Due to budget constraints, we can provide evaluators with one H100 80GB GPU instance on a public cloud for approximately 24 hours.

\subsubsection{Software dependencies}

(1) \texttt{conda}, Python~3.10, \texttt{uv} (managed by
\texttt{rebuild\_env.sh}), \texttt{matplotlib}; HuggingFace access token for Llama-3 gated models.
(2) Python~3.10, \texttt{tqdm}, \texttt{pandas}, \texttt{matplotlib}
(installed via \texttt{pip install -r requirements.txt}).

\subsubsection{Datasets}

(1) All datasets except ShareGPT are downloaded automatically via
HuggingFace \texttt{datasets}. ShareGPT requires a one-time download
(see \S\ref{sec:ae-install}).
(2) Pre-profiled JSONL acceptance data for Llama-3.1-8B-Instruct on
H100 is included in \texttt{simulator/data/llama3.1-8B/} (gsm8k,
instructcoder, sharegpt, cnn).

%%%%%%%%%%%%%%%%%%%%%%%%%%%%%%%%%%%%%%%%%%%%%%%%%%%%%%%%%%%%%%%%%%%%%
\subsection{Installation}
\label{sec:ae-install}

\textbf{(1) vLLM suite: ShareGPT} (one-time):
\begin{verbatim}
huggingface-cli download  \
anon8231489123/ShareGPT_Vicuna_unfiltered \
    ShareGPT_V3_unfiltered_cleaned_split.json \
    --repo-type dataset --local-dir /path/to/data/
export SHAREGPT_PATH=/path/to/dataset
\end{verbatim}

\textbf{(1) vLLM suite: environment} (repeat per branch):
\begin{verbatim}
git checkout <branch>
ENV_DIR=/path/to/envs bash scripts/rebuild_env.sh
conda activate /path/to/envs
\end{verbatim}

\textbf{(2) Simulator}:
\begin{verbatim}
cd simulator/
conda create -n specbench_simulator python=3.10 -y
conda activate specbench_simulator
pip install -r requirements.txt
\end{verbatim}

%%%%%%%%%%%%%%%%%%%%%%%%%%%%%%%%%%%%%%%%%%%%%%%%%%%%%%%%%%%%%%%%%%%%%
\subsection{Experiment workflow}

\textbf{(1) vLLM profiling.}
Set \texttt{CUDA\_VISIBLE\_DEVICES} at the top of the target script,
then from \texttt{scripts/}:
\begin{verbatim}
export SHAREGPT_PATH=/path/to/sharegpt.json
bash run-l3-8b.sh    # 1 GPU
\end{verbatim}
Each script runs warmup, all datasets $\times$ methods, and calls
\texttt{vis\_speedup.py} to produce PDF figures in
\texttt{results/run\_<timestamp>/figures/}.

\textbf{(2) Simulator.}
\begin{verbatim}
cd simulator/
./simulate_and_plot.sh
\end{verbatim}
This runs EAGLE-3, N-gram, and combined proposers across all four
datasets and saves CSV results to \texttt{results/} and PDF figures to
\texttt{figures/llama3.1-8B/}.

%%%%%%%%%%%%%%%%%%%%%%%%%%%%%%%%%%%%%%%%%%%%%%%%%%%%%%%%%%%%%%%%%%%%%
\subsection{Evaluation and expected result}

\textbf{(1) vLLM profiling.}
After running \texttt{run-l3-8b.sh}, we expect the same end-to-end performance results to be reproduced for Figure~1\,(a)--(d).
% Expected behavior on H100:
% \begin{itemize}
%   \item \textbf{Llama-3.1-8B} (\texttt{run-l3-8b.sh}, reproduces
%     Figure~1\,(a)--(d)): EAGLE-3 achieves the highest speedup
%     ($\approx$1.5--2.2$\times$ at batch\,=\,1), followed by EAGLE
%     ($\approx$1.7$\times$ on GSM8K at batch\,=\,1, dropping to
%     $\approx$1.2$\times$ at batch\,=\,128). N-gram outperforms
%     EAGLE/EAGLE-3 on InstructCoder due to high prompt-output overlap.
%   \item \textbf{Llama-3-70B} (\texttt{run-l3-70b.sh}, requires 4$\times$H100):
%     Draft-model outperforms EAGLE on most settings
%     ($\approx$1.7--2.0$\times$ at batch\,=\,1).
%   \item \textbf{Qwen3-8B} (\texttt{run-q3-8b.sh}, reasoning mode,
%     $\approx$48\,h): EAGLE-3 achieves $\approx$1.6--1.8$\times$ speedup
%     on GPQA-Main and AIME22-24.
% \end{itemize}
% Results may vary by $\pm$5--10\% across GPU generations. Check \texttt{overview.log} for per-run success/failure.

\textbf{(2) Simulator.} We expect the same results to be reproduced for Figure 9 and Figure 11.

%%%%%%%%%%%%%%%%%%%%%%%%%%%%%%%%%%%%%%%%%%%%%%%%%%%%%%%%%%%%%%%%%%%%%
\subsection{Experiment customization}

\textbf{(1) vLLM profiling.}
Set \texttt{num\_reqs=5} and \texttt{batch\_sizes="1 16"} for a
$<$10\,min smoke test (use a non-reasoning script).
Individual datasets and methods can be toggled by editing the
\texttt{for} loops.

\textbf{(2) Simulator.}
Individual proposers, datasets, prediction methods, and batch sizes
can be configured via CLI flags to \texttt{main.py}. Refer to the README on GitHub for more information.

%%%%%%%%%%%%%%%%%%%%%%%%%%%%%%%%%%%%%%%%%%%%%%%%%%%%%%%%%%%%%%%%%%%%%
\subsection{Notes}

For profiling, each script run writes to a fresh timestamped directory. Each branch requires its own virtual environment.

%% <<< end appendix.tex


% \section{Electronic Submission}
% \label{submission}

% Submission to MLSys 2025 will be entirely electronic, via a web site
% (not email). Information about the submission process and \LaTeX\ templates
% are available on the conference web site at:
% \begin{center}
% \textbf{\texttt{https://www.mlsys.org/}}
% \end{center}

% The guidelines below will be enforced for initial submissions and
% camera-ready copies. Here is a brief summary:
% \begin{itemize}
% \item Submissions must be in PDF\@.
% \item The maximum paper length is \textbf{10 pages} excluding references and appendices.
%     (pages 11 and on must contain only references).
% \item \textbf{Do not include author information or acknowledgements} in your
%     initial submission.
% \item Your paper should be in \textbf{10 point Times font}.
% \item Make sure your PDF file only uses Type-1 fonts.
% \item Place figure captions \emph{under} the figure (and omit titles from inside
%     the graphic file itself). Place table captions \emph{over} the table.
% \item References must include page numbers whenever possible and be as complete
%     as possible. Place multiple citations in chronological order.
% \item Do not alter the style template; in particular, do not compress the paper
%     format by reducing the vertical spaces.
% \item Keep your abstract brief and self-contained, one paragraph and roughly
%     4--6 sentences. Gross violations will require correction at the
%     camera-ready phase. The title should have content words capitalized.
% \item Submissions that significantly diverge from the specified format 
%     will be rejected without review.
% \end{itemize}

% \subsection{Submitting Papers}

% \textbf{Paper Deadline:} The deadline for paper submission that is
% advertised on the conference website is strict. If your full,
% anonymized, submission does not reach us on time, it will not be
% considered for publication. There is no separate abstract submission.

% \textbf{Anonymous Submission:} MLSys uses double-blind review: no identifying
% author information may appear on the title page or in the paper
% itself. Section~\ref{author info} gives further details.

% \textbf{Simultaneous Submission:} MLSys will not accept any paper which,
% at the time of submission, is under review for another conference or
% has already been published. This policy also applies to papers that
% overlap substantially in technical content with conference papers
% under review or previously published. MLSys submissions must not be
% submitted to other conferences during MLSys's review period. Authors
% may submit to MLSys substantially different versions of journal papers
% that are currently under review by the journal, but not yet accepted
% at the time of submission. Informal publications, such as technical
% reports or papers in workshop proceedings which do not appear in
% print, do not fall under these restrictions.

% \medskip

% Authors must provide their manuscripts in \textbf{PDF} format.
% Furthermore, please make sure that files contain only embedded Type-1 fonts
% (e.g.,~using the program \texttt{pdffonts} in linux or using
% File/DocumentProperties/Fonts in Acrobat). Other fonts (like Type-3)
% might come from graphics files imported into the document.

% Authors using \textbf{Word} must convert their document to PDF\@. Most
% of the latest versions of Word have the facility to do this
% automatically. Submissions will not be accepted in Word format or any
% format other than PDF\@.

% Those who use \textbf{\LaTeX} should avoid including Type-3 fonts.
% Those using \texttt{latex} and \texttt{dvips} may need the following
% two commands:

% {\footnotesize
% \begin{verbatim}
% dvips -Ppdf -tletter -G0 -o paper.ps paper.dvi
% ps2pdf paper.ps
% \end{verbatim}}
% It is a zero following the ``-G'', which tells dvips to use
% the config.pdf file. Newer \TeX\ distributions don't always need this
% option.

% Using \texttt{pdflatex} rather than \texttt{latex}, often gives better
% results. This program avoids the Type-3 font problem, and supports more
% advanced features in the \texttt{microtype} package.

% \textbf{Graphics files} should be a reasonable size, and included from
% an appropriate format. Use vector formats (.eps/.pdf) for plots,
% lossless bitmap formats (.png) for raster graphics with sharp lines, and
% jpeg for photo-like images.

% The style file uses the \texttt{hyperref} package to make clickable
% links in documents. If this causes problems for you, add
% \texttt{nohyperref} as one of the options to the \texttt{mlsys2025}
% usepackage statement.


% \subsection{Submitting Final Camera-Ready Copy}

% The final versions of papers accepted for publication should follow the
% same format and naming convention as initial submissions, except that
% author information (names and affiliations) should be given. See
% Section~\ref{final author} for formatting instructions.

% The footnote, ``Preliminary work. Under review by the MLSys. Do not distribute.'' 
% must be modified to ``\textit{Proceedings of the
% $\mathit{8}^{th}$ MLSys Conference},
% Santa Clara, CA, USA, 2025.
% Copyright 2025 by the author(s).''

% For those using the \textbf{\LaTeX} style file, this change (and others) is
% handled automatically by simply changing
% $\mathtt{\backslash usepackage\{mlsys2025\}}$ to
% $$\mathtt{\backslash usepackage[accepted]\{mlsys2025\}}$$
% Authors using \textbf{Word} must edit the
% footnote on the first page of the document themselves.

% Camera-ready copies should have the title of the paper as running head
% on each page except the first one. The running title consists of a
% single line centered above a horizontal rule which is $1$~point thick.
% The running head should be centered, bold and in $9$~point type. The
% rule should be $10$~points above the main text. For those using the
% \textbf{\LaTeX} style file, the original title is automatically set as running
% head using the \texttt{fancyhdr} package which is included in the MLSys
% 2025 style file package. In case that the original title exceeds the
% size restrictions, a shorter form can be supplied by using

% \verb|\mlsystitlerunning{...}|

% just before $\mathtt{\backslash begin\{document\}}$.
% Authors using \textbf{Word} must edit the header of the document themselves.

% \section{Format of the Paper}

% All submissions must follow the specified format.

% \subsection{Length and Dimensions}

% Papers must not exceed ten (10) pages, including all figures, tables,
% but excluding acknowledgements, references and appendices.
% Acknowledgements should be limited to grants and people who contributed to the paper.
% Any submission that exceeds
% this page limit, or that diverges significantly from the specified format,
% will be rejected without review.

% The text of the paper should be formatted in two columns, with an
% overall width of 6.75~inches, height of 9.0~inches, and 0.25~inches
% between the columns. The left margin should be 0.75~inches and the top
% margin 1.0~inch (2.54~cm). The right and bottom margins will depend on
% whether you print on US letter or A4 paper, but all final versions
% must be produced for US letter size.

% The paper body should be set in 10~point type with a vertical spacing
% of 11~points. Please use Times typeface throughout the text.

% \subsection{Title}

% The paper title should be set in 14~point bold type and centered
% between two horizontal rules that are 1~point thick, with 1.0~inch
% between the top rule and the top edge of the page. Capitalize the
% first letter of content words and put the rest of the title in lower
% case.

% \subsection{Author Information for Submission}
% \label{author info}

% MLSys uses double-blind review, so author information must not appear. If
% you are using \LaTeX\/ and the \texttt{mlsys2025.styappen} file, use
% \verb+\mlsysauthor{...}+ to specify authors and \verb+\mlsysaffiliation{...}+ to specify affiliations. (Read the TeX code used to produce this document for an example usage.) The author information
% will not be printed unless \texttt{accepted} is passed as an argument to the
% style file.
% Submissions that include the author information will not
% be reviewed.

% \subsubsection{Self-Citations}

% If you are citing published papers for which you are an author, refer
% to yourself in the third person. In particular, do not use phrases
% that reveal your identity (e.g., ``in previous work \cite{langley00}, we
% have shown \ldots'').

% Do not anonymize citations in the reference section. The only exception are manuscripts that are
% not yet published (e.g., under submission). If you choose to refer to
% such unpublished manuscripts \cite{anonymous}, anonymized copies have
% to be submitted
% as Supplementary Material via CMT\@. However, keep in mind that an MLSys
% paper should be self contained and should contain sufficient detail
% for the reviewers to evaluate the work. In particular, reviewers are
% not required to look at the Supplementary Material when writing their
% review.

% \subsubsection{Camera-Ready Author Information}
% \label{final author}

% If a paper is accepted, a final camera-ready copy must be prepared.
% %
% For camera-ready papers, author information should start 0.3~inches below the
% bottom rule surrounding the title. The authors' names should appear in 10~point
% bold type, in a row, separated by white space, and centered. Author names should
% not be broken across lines. Unbolded superscripted numbers, starting 1, should
% be used to refer to affiliations.

% Affiliations should be numbered in the order of appearance. A single footnote
% block of text should be used to list all the affiliations.

% Each distinct affiliations should be listed once. If an author has multiple
% affiliations, multiple superscripts should be placed after the name, separated
% by thin spaces. If the authors would like to highlight equal contribution by
% multiple first authors, those authors should have an asterisk placed after their
% name in superscript, and the term ``\textsuperscript{*}Equal contribution"
% should be placed in the footnote block ahead of the list of affiliations. A
% list of corresponding authors and their emails (in the format Full Name
% \textless{}email@domain.com\textgreater{}) can follow the list of affiliations.
% Ideally only one or two names should be listed.

% A sample file with author names is included in the mlsys2025 style file
% package. Turn on the \texttt{[accepted]} option to the stylefile to
% see the names rendered. All of the guidelines above are implemented
% by the \LaTeX\ style file.

% \subsection{Abstract}

% The paper abstract should begin in the left column, 0.4~inches below the final
% address. The heading `Abstract' should be centered, bold, and in 11~point type.
% The abstract body should use 10~point type, with a vertical spacing of
% 11~points, and should be indented 0.25~inches more than normal on left-hand and
% right-hand margins. Insert 0.4~inches of blank space after the body. Keep your
% abstract brief and self-contained, limiting it to one paragraph and roughly 4--6
% sentences. Gross violations will require correction at the camera-ready phase.

% \subsection{Partitioning the Text}

% You should organize your paper into sections and paragraphs to help
% readers place a structure on the material and understand its
% contributions.

% \subsubsection{Sections and Subsections}

% Section headings should be numbered, flush left, and set in 11~pt bold
% type with the content words capitalized. Leave 0.25~inches of space
% before the heading and 0.15~inches after the heading.

% Similarly, subsection headings should be numbered, flush left, and set
% in 10~pt bold type with the content words capitalized. Leave
% 0.2~inches of space before the heading and 0.13~inches afterward.

% Finally, subsubsection headings should be numbered, flush left, and
% set in 10~pt small caps with the content words capitalized. Leave
% 0.18~inches of space before the heading and 0.1~inches after the
% heading.

% Please use no more than three levels of headings.

% \subsubsection{Paragraphs and Footnotes}

% Within each section or subsection, you should further partition the
% paper into paragraphs. Do not indent the first line of a given
% paragraph, but insert a blank line between succeeding ones.

% You can use footnotes\footnote{Footnotes
% should be complete sentences.} to provide readers with additional
% information about a topic without interrupting the flow of the paper.
% Indicate footnotes with a number in the text where the point is most
% relevant. Place the footnote in 9~point type at the bottom of the
% column in which it appears. Precede the first footnote in a column
% with a horizontal rule of 0.8~inches.\footnote{Multiple footnotes can
% appear in each column, in the same order as they appear in the text,
% but spread them across columns and pages if possible.}

% \subsection{Figures}

% You may want to include figures in the paper to illustrate
% your approach and results. Such artwork should be centered,
% legible, and separated from the text. Lines should be dark and at
% least 0.5~points thick for purposes of reproduction, and text should
% not appear on a gray background.

% Label all distinct components of each figure. If the figure takes the
% form of a graph, then give a name for each axis and include a legend
% that briefly describes each curve. Do not include a title inside the
% figure; instead, the caption should serve this function.

% Number figures sequentially, placing the figure number and caption
% \emph{after} the graphics, with at least 0.1~inches of space before
% the caption and 0.1~inches after it. The figure caption should be set in
% 9~point type and centered unless it runs two or more lines, in which
% case it should be flush left. You may float figures to the top or
% bottom of a column, and you may set wide figures across both columns
% (use the environment \texttt{figure*} in \LaTeX). Always place
% two-column figures at the top or bottom of the page.

% \subsection{Algorithms}

% If you are using \LaTeX, please use the ``algorithm'' and ``algorithmic''
% environments to format pseudocode. These require
% the corresponding stylefiles, algorithm.sty and
% algorithmic.sty, which are supplied with this package.
% Algorithm~\ref{alg:example} shows an example.

% \begin{algorithm}[tb]
%    \caption{Bubble Sort}
%    \label{alg:example}
% \begin{algorithmic}
%    \STATE {\bfseries Input:} data $x_i$, size $m$
%    \REPEAT
%    \STATE Initialize $noChange = true$.
%    \FOR{$i=1$ {\bfseries to} $m-1$}
%    \IF{$x_i > x_{i+1}$}
%    \STATE Swap $x_i$ and $x_{i+1}$
%    \STATE $noChange = false$
%    \ENDIF
%    \ENDFOR
%    \UNTIL{$noChange$ is $true$}
% \end{algorithmic}
% \end{algorithm}

% \subsection{Tables}

% You may also want to include tables that summarize material. Like
% figures, these should be centered, legible, and numbered consecutively.
% However, place the title \emph{above} the table with at least
% 0.1~inches of space before the title and the same after it, as in
% Table~\ref{sample-table}. The table title should be set in 9~point
% type and centered unless it runs two or more lines, in which case it
% should be flush left.

% % Note use of \abovespace and \belowspace to get reasonable spacing
% % above and below tabular lines.

% \begin{table}[t]
% \caption{Classification accuracies for naive Bayes and flexible
% Bayes on various data sets.}
% \label{sample-table}
% \vskip 0.15in
% \begin{center}
% \begin{small}
% \begin{sc}
% \begin{tabular}{lcccr}
% \toprule
% Data set & Naive & Flexible & Better? \\
% \midrule
% Breast    & 95.9$\pm$ 0.2& 96.7$\pm$ 0.2& $\surd$ \\
% Cleveland & 83.3$\pm$ 0.6& 80.0$\pm$ 0.6& $\times$\\
% Glass2    & 61.9$\pm$ 1.4& 83.8$\pm$ 0.7& $\surd$ \\
% Credit    & 74.8$\pm$ 0.5& 78.3$\pm$ 0.6&         \\
% Horse     & 73.3$\pm$ 0.9& 69.7$\pm$ 1.0& $\times$\\
% Meta      & 67.1$\pm$ 0.6& 76.5$\pm$ 0.5& $\surd$ \\
% Pima      & 75.1$\pm$ 0.6& 73.9$\pm$ 0.5&         \\
% Vehicle   & 44.9$\pm$ 0.6& 61.5$\pm$ 0.4& $\surd$ \\
% \bottomrule
% \end{tabular}
% \end{sc}
% \end{small}
% \end{center}
% \vskip -0.1in
% \end{table}

% Tables contain textual material, whereas figures contain graphical material.
% Specify the contents of each row and column in the table's topmost
% row. Again, you may float tables to a column's top or bottom, and set
% wide tables across both columns. Place two-column tables at the
% top or bottom of the page.

% \subsection{Citations and References}

% Please use APA reference format regardless of your formatter
% or word processor. If you rely on the \LaTeX\/ bibliographic
% facility, use \texttt{natbib.sty} and \texttt{mlsys2025.bst}
% included in the style-file package to obtain this format.

% Citations within the text should include the authors' last names and
% year. If the authors' names are included in the sentence, place only
% the year in parentheses, for example when referencing Arthur Samuel's
% pioneering work \yrcite{Samuel59}. Otherwise place the entire
% reference in parentheses with the authors and year separated by a
% comma \cite{Samuel59}. List multiple references separated by
% semicolons \cite{kearns89,Samuel59,mitchell80}. Use the `et~al.'
% construct only for citations with three or more authors or after
% listing all authors to a publication in an earlier reference \cite{MachineLearningI}.

% Authors should cite their own work in the third person
% in the initial version of their paper submitted for blind review.
% Please refer to Section~\ref{author info} for detailed instructions on how to
% cite your own papers.

% Use an unnumbered first-level section heading for the references, and use a
% hanging indent style, with the first line of the reference flush against the
% left margin and subsequent lines indented by 10 points. The references at the
% end of this document give examples for journal articles \cite{Samuel59},
% conference publications \cite{langley00}, book chapters \cite{Newell81}, books
% \cite{DudaHart2nd}, edited volumes \cite{MachineLearningI}, technical reports
% \cite{mitchell80}, and dissertations \cite{kearns89}.

% Alphabetize references by the surnames of the first authors, with
% single author entries preceding multiple author entries. Order
% references for the same authors by year of publication, with the
% earliest first. Make sure that each reference includes all relevant
% information (e.g., page numbers).

% Please put some effort into making references complete, presentable, and
% consistent. If using bibtex, please protect capital letters of names and
% abbreviations in titles, for example, use \{B\}ayesian or \{L\}ipschitz
% in your .bib file.

% \subsection{Software and Data}

% We strongly encourage the publication of software and data with the
% camera-ready version of the paper whenever appropriate. This can be
% done by including a URL in the camera-ready copy. However, do not
% include URLs that reveal your institution or identity in your
% submission for review. Instead, provide an anonymous URL or upload
% the material as ``Supplementary Material'' into the CMT reviewing
% system. Note that reviewers are not required to look at this material
% when writing their review.

% % Acknowledgements should only appear in the accepted version.
% \section*{Acknowledgements}

% \textbf{Do not} include acknowledgements in the initial version of
% the paper submitted for blind review.

% If a paper is accepted, the final camera-ready version can (and
% probably should) include acknowledgements. In this case, please
% place such acknowledgements in an unnumbered section at the
% end of the paper. Typically, this will include thanks to reviewers
% who gave useful comments, to colleagues who contributed to the ideas,
% and to funding agencies and corporate sponsors that provided financial
% support.


% % In the unusual situation where you want a paper to appear in the
% % references without citing it in the main text, use \nocite
% \nocite{langley00}




% %%%%%%%%%%%%%%%%%%%%%%%%%%%%%%%%%%%%%%%%%%%%%%%%%%%%%%%%%%%%%%%%%%%%%%%%%%%%%%%
% %%%%%%%%%%%%%%%%%%%%%%%%%%%%%%%%%%%%%%%%%%%%%%%%%%%%%%%%%%%%%%%%%%%%%%%%%%%%%%%
% % SUPPLEMENTAL CONTENT AS APPENDIX AFTER REFERENCES
% %%%%%%%%%%%%%%%%%%%%%%%%%%%%%%%%%%%%%%%%%%%%%%%%%%%%%%%%%%%%%%%%%%%%%%%%%%%%%%%
% %%%%%%%%%%%%%%%%%%%%%%%%%%%%%%%%%%%%%%%%%%%%%%%%%%%%%%%%%%%%%%%%%%%%%%%%%%%%%%%
% \appendix
% \section{Please add supplemental material as appendix here}
% %
% Put anything that you might normally include after the references as an appendix here, {\it not in a separate supplementary file}. Upload your final camera-ready as a single pdf, including all appendices.

% %%%%%%%%%%%%%%%%%%%%%%%%%%%%%%%%%%%%%%%%%%%%%%%%%%%%%%%%%%%%%%%%%%%%%%%%%%%%%%%
% %%%%%%%%%%%%%%%%%%%%%%%%%%%%%%%%%%%%%%%%%%%%%%%%%%%%%%%%%%%%%%%%%%%%%%%%%%%%%%%





\end{document}


% This document was modified from the file originally made available by
% Pat Langley and Andrea Danyluk for ICML-2K. This version was created
% by Iain Murray in 2018. It was modified from a version from Dan Roy in
% 2017, which was based on a version from Lise Getoor and Tobias
% Scheffer, which was slightly modified from the 2010 version by
% Thorsten Joachims & Johannes Fuernkranz, slightly modified from the
% 2009 version by Kiri Wagstaff and Sam Roweis's 2008 version, which is
% slightly modified from Prasad Tadepalli's 2007 version which is a
% lightly changed version of the previous year's version by Andrew
% Moore, which was in turn edited from those of Kristian Kersting and
% Codrina Lauth. Alex Smola contributed to the algorithmic style files.
